% !TEX TS-program = lualatex
% Documentation of tkz-elements v5.04c
% Copyright (c) 2023–2025 Alain Matthes
% SPDX-License-Identifier: LPPL-1.3c
% Maintainer: Alain Matthes

\documentclass[
  DIV=16,
  fontsize=10,
  index=totoc,
  twoside,
  cadre,
  headings=small
]{tkz-doc}
%\overfullrule=5pt
\gdef\tkznameofpack{tkz-elements}
\gdef\tkzversionofpack{5.05c}
\gdef\tkzdateofpack{\today}
\gdef\tkznameofdoc{tkz-elements.pdf}
\let\tkzversionofdoc\tkzversionofpack
\let\tkzdateofdoc\tkzdateofpack
\gdef\tkzauthorofpack{Alain Matthes}
\gdef\tkzadressofauthor{}
\gdef\tkznamecollection{AlterMundus}
\gdef\tkzurlauthor{http://altermundus.fr}
\gdef\tkzengine{lualatex}
\nameoffile{\tkznameofpack}
\let\tkzurlauthorcom\tkzurlauthor
% -- Packages ---------------------------------------------------
\usepackage[dvipsnames,svgnames]{xcolor}
% my packages
\usepackage{tkz-base}
\usetikzlibrary{mindmap,shapes.multipart}
\usepackage[mini]{tkz-euclide}
\usepackage{tkz-elements}
\usepackage{tkzexample}
\usepackage{pgfornament}

% mdf warning silence (optionnel)
\makeatletter
    \renewcommand*\mdf@PackageWarning[1]{}
     \renewcommand*\mdf@PackageInfo[1]{}
\makeatother

% for metapost
\usepackage{luamplib}
\everymplib{input mpcolornames; beginfig(1);}
\everyendmplib{endfig;}

% for index
\usepackage{scrindex}
\makeindex


\usepackage[colorlinks,pdfencoding=auto, psdextra]{hyperref}
\hypersetup{
    linkcolor=black,
    citecolor=Green,
    filecolor=Mulberry,
    urlcolor=orange,
    menucolor=black,
    runcolor=Mulberry,
    linkbordercolor=black,
    citebordercolor=Green,
    filebordercolor=Mulberry,
    urlbordercolor=NavyBlue,
    menubordercolor=black,
    runbordercolor=Mulberry,
    pdfsubject={Euclidean Geometry},
    pdfauthor={\tkzauthorofpack},
    pdftitle={\tkznameofpack},
    pdfcreator={\tkzengine}
}

% ...then bookmark
\usepackage{bookmark}

% fonts
%<--------------------------------------------------------------------------->
\usepackage{fontspec}
\setmainfont{Linux Libertine O}
\usepackage{unicode-math}
\usepackage{concmath-otf}
\let\rmfamily\ttfamily
%<--------------------------------------------------------------------------->
% load other packages
\usepackage{pifont}
\newcommand{\cmark}{\colorbox{green}{\textcolor{white}{\ding{51}}}}
\newcommand{\xmark}{\colorbox{red}{\textcolor{white}{\ding{55}}}}
\usepackage{multicol}
\usepackage[english]{babel}
\usepackage[normalem]{ulem}
\usepackage{booktabs}
\usepackage{array}
\usepackage{cellspace}
\usepackage{shortvrb}
\usepackage{fancyvrb}
\usepackage{enumitem}
\usepackage{moreenum}
\usepackage{textcomp}
\usepackage{mathtools}
\usepackage{amsthm}
\usepackage{thmtools}

\usepackage[most]{tcolorbox}
\newtcolorbox{mybox}{
  enhanced,
  boxrule=0pt,frame hidden,
  borderline west={3pt}{0pt}{cyan!75!black},
  colback=lightgray!10!white,
  sharp corners
}
% end load packages
%<--------------------------------------------------------------------------->
% Indexation style
\makeatletter
\renewenvironment{theindex}
   {\section*{\indexname}\begin{multicols}{2}%
    \@mkboth{\MakeUppercase\indexname}%
  {\MakeUppercase\indexname}%
    \thispagestyle{plain}\parindent\z@
    \parskip\z@ \@plus .3\p@\relax
    \columnseprule \z@
    \columnsep 35\p@
    \let\item\@idxitem}
   {\end{multicols}}
\makeatother
% end indexation style
%<--------------------------------------------------------------------------->
% Indexation
\newcommand{\tkzAttr}[2]{\textbf{\texttt{#2}}\index{#1_3@\texttt{#1: Attributes}!\texttt{#2}}}
\newcommand{\tkzVar}[2]{\textbf{\texttt{#2}}\index{#1_3@\texttt{#1: Reserved variable}!\texttt{#2}}}
\newcommand{\tkzMeth}[2]{\textbf{\texttt{#2}}\index{#1_3@\texttt{#1: Methods}!\texttt{#2}}}
\newcommand{\tkzMeta}[2]{\textbf{\texttt{#2}}\index{#1_3@\texttt{#1: Metamethods}!\texttt{#2}}}
\newcommand{\tkzFct}[2]{\textbf{\texttt{#2}}\index{#1_3@\texttt{#1: Functions}!\texttt{#2}}}
\newcommand{\tkzMod}[1]{\textbf{\texttt{#1}}\index{Modules!#1@\texttt{#1}}}
\newcommand{\tkzClass}[1]{\textbf{\texttt{#1}}\index{Classes!#1@\texttt{#1}}}
\newcommand{\tkzNameObj}[1]{\textbf{\texttt{#1}}\index{Objects!#1@\texttt{#1}}}
\newcommand{\tkzMacro}[2]{\textbf{\texttt{\textbackslash#2}}\index{#1_3@\texttt{#1: Macro}!\texttt{#2}}}
\newcommand{\tkzEngine}[1]{\textbf{\texttt{#1}}\index{Engine !#1@\texttt{#1}}}
\newcommand{\tkzPopt}[2]{\texttt{#2}\index{#1_3@\texttt{#1: options}!\texttt{#2}}}
\newcommand{\tkzEnv}[2]{\textbf{\texttt{#2}}\index{#1_3@\texttt{#1: Environments}!\texttt{#2}}}
\newcommand{\tkzConst}[1]{\textbf{\texttt{#1}}\index{Constants!#1@\texttt{#1}}}
\newcommand{\tkzSysVar}[1]{\textbf{\texttt{#1}}\index{Variables system!#1@\texttt{#1}}}
% Keyword for function options or parameters
\newcommand{\tkzKey}[1]{\texttt{#1}}

% Optional argument notation, e.g., (pt)
\newcommand{\tkzOptArg}[1]{\ensuremath{(#1)}}
% Keyword used as an optional key in a function call (e.g., known, near)
%\newcommand{\tkzKey}[1]{\texttt{#1}\index{#1@\texttt{#1}}}

% end indexation
%<--------------------------------------------------------------------------->
% decoration
\newcommand{\lefthand}{\ding{43}}
\newcommand{\tkzimpbf}[1]{\texttt{tkzPopt#1}}
\newcommand{\tkzRHand}{\textcolor{red}{\lefthand}}
\newcommand{\smallbf}[1]{\textbf{\small #1}}
\newcommand{\parallelsum}{\mathbin{\!/\mkern-5mu/\!}}
\newcommand\LUALATEX{Lua\kern-.01em\LaTeX}%
% end decoration
%<--------------------------------------------------------------------------->
% settings styles
\tkzSetUpColors[background=white,text=black]
\tkzSetUpPoint[size=2,color=teal,fill=teal!10]
\tkzSetUpLine[ultra thin,color=teal]
\tkzSetUpCompass[color=orange,ultra thin,/tkzcompass/delta=10]
\tikzset{label style/.append style={below,color=teal,font=\scriptsize}}
\tikzset{new/.style={color=orange,ultra thin}}
\tikzset{step 1/.style={color=cyan,ultra thin}}
\tikzset{step 2/.style={color=purple,ultra thin}}
\def\tkzar{\hspace{1em}-->\hspace{1em}}
% end settings styles
%<--------------------------------------------------------------------------->
% misc
\makeatletter\let\percentchar\@percentchar\makeatother
\def\luaveclen#1#2{%
\directlua{tex.print(string.format(
'\percentchar.5f',math.sqrt((#1)*(#1)+(#2)*(#2))))
}}
\let\pmpn\pgfmathprintnumber
\let\code\tkzname
% end misc
%<--------------------------------------------------------------------------->
% Generate a new TOC
%----------------------------------------------------------
% Mini-TOC "Overview" (stoc) with styled parts
%----------------------------------------------------------
%----------------------------------------------------------
% Mini-TOC "Overview" (stoc) : now includes \part + \section
%----------------------------------------------------------
%==========================================================
% Mini-TOC "Overview" (stoc)
%==========================================================

\DeclareNewTOC[
  type=overview,
  listname={Overview}
]{stoc}

%----------------------------------------------------------
% Save original KOMA macros
%----------------------------------------------------------

\NewCommandCopy{\AddPartTocEntry}{\addparttocentry}
\NewCommandCopy{\AddSectionTocEntry}{\addsectiontocentry}

%----------------------------------------------------------
% Mirror \part into stoc (+ blank line after)
%----------------------------------------------------------

\renewcommand{\addparttocentry}[2]{%
  \AddPartTocEntry{#1}{#2}%
  \addxcontentsline{stoc}{part}[{#1}]{#2}%
  \addtocontents{stoc}{\protect\addvspace{\baselineskip}}%
}

%----------------------------------------------------------
% Mirror \section into stoc
%----------------------------------------------------------

\renewcommand{\addsectiontocentry}[2]{%
  \AddSectionTocEntry{#1}{#2}%
  \addxcontentsline{stoc}{section}[{#1}]{#2}%
}

%----------------------------------------------------------
% Style of the Overview (stoc)
%----------------------------------------------------------

\BeforeStartingTOC[stoc]{%

  % PART entries — code style
  \DeclareTOCStyleEntry[
    beforeskip=.6\baselineskip,
    entryformat=\ttfamily\bfseries,
    pagenumberformat=\ttfamily\bfseries,
    linefill=\hfill
  ]{part}{part}

  % SECTION entries — compact
  \DeclareTOCStyleEntry[
    beforeskip=0pt
  ]{section}{section}
}
%==========================================================
%--------------------------------------------------------------------------------------------------------------------

\hbadness=10000
%<--------------------------------------------------------------------------->
\AtBeginDocument{\MakeShortVerb{\|}}

\begin{document}

\parindent=0pt
\tkzTitleFrame{tkz-elements \tkzversionofpack\\Euclidean Geometry}
\clearpage


\defoffile{
	This document gathers notes on \tkzNamePack{\tkznameofpack}, a comprehensive \tkzname{Lua} library designed to handle all necessary computations for constructing objects in Euclidean geometry. The document must be compiled with Lua\LATEX.\\

	With \tkzNamePack{tkz-elements}, all definitions and calculations are carried out exclusively in Lua, following an object-oriented programming model. Core geometric entities such as points, lines, triangles, and conics are implemented as object classes.\\
	Once the computations are complete, graphical rendering is typically handled using \tkzNamePack{tkz-euclide}, although you may also use \TIKZ{} directly if desired.\\

	I discovered Lua and object-oriented programming while developing this package, so it’s possible that some parts could be improved. If you'd like to contribute or offer suggestions, feel free to contact me by email.
}

\presentation

\vspace*{4em}

\lefthand\ Acknowledgements: I received much valuable advice, remarks, and corrections  from \\ \tkzimp{Nicolas Kisselhoff}, \tkzimp{David Carlisle}, \tkzimp{Roberto Giacomelli} and \tkzimp{Qrrbrbirlbel}.\\
Special thanks to \tkzimp{Wolfgang Büchel} for his invaluable contribution in correcting the examples.

\vspace*{1em}
\lefthand\ I would also like to extend my gratitude to \tkzimp{Eric Weisstein}, creator of
\href{http://mathworld.wolfram.com/about/author.html}{MathWorld}.

\vspace*{1em}
\lefthand\ You can find some examples on my site and a french documentation:
\href{http://altermundus.fr}{altermundus.fr}.

\vfill
Please report typos or any other comments to this documentation to: \href{mailto:al.ma@mac.com}{\textcolor{blue}{Alain Matthes}}.

This file can be redistributed and/or modified under the terms of the \LaTeX{}
Project Public License Distributed from \href{http://www.ctan.org/}{CTAN}\  archives.

\clearpage
\listofoverviews

\clearpage
\section*{Attributes Overview}
\begin{tabular}{ll}
\textbf{Class} & \textbf{Attributes} \\
\hline
point         & \hyperref[point:attributes]{see table}        \\
line          & \hyperref[line:attributes]{see table}         \\
triangle      & \hyperref[triangle:attributes]{see table}     \\
circle        & \hyperref[circle:attributes]{see table}       \\
occs          & \hyperref[occs:attributes]{see table}         \\
conic         & \hyperref[conic:attributes]{see table}        \\
quadrilateral & \hyperref[quadrilateral:attributes]{see table}\\
square        & \hyperref[square:attributes]{see table}       \\
rectangle     & \hyperref[rectangle:attributes]{see table}    \\
parallelogram & \hyperref[parallelogram:attributes]{see table}\\
regular       & \hyperref[regular:attributes]{see table}      \\
vector        & \hyperref[vector:attributes]{see table}       \\
matrix        & \hyperref[matrix:attributes]{see table}       \\
angle         & \hyperref[angle:attributes]{see table}        \\
list\_point   & \hyperref[list_point:attributes]{see table}   \\
\end{tabular}
\section*{Methods Overview}
\begin{tabular}{ll}
\textbf{Class} & \textbf{Methods} \\
\hline
point         & \hyperref[point:methods]{see table}         \\
line          & \hyperref[line:methods part 1]{see table 1} \\
line          & \hyperref[line:methods part 2]{see table 2} \\
triangle      & \hyperref[triangle:methods]{see table}      \\
circle        & \hyperref[circle:methods]{see table}        \\
circle        & \hyperref[circle:methods 2]{see table 2}    \\
occs          & \hyperref[occs:methods]{see table}          \\
conic         & \hyperref[conic:methods]{see table}         \\
path          & \hyperref[path:methods]{see table}          \\
quadrilateral & \hyperref[quadrilateral:methods]{see table} \\
square        & \hyperref[square:methods]{see table}        \\
rectangle     & \hyperref[rectangle:methods]{see table}     \\
parallelogram & \hyperref[parallelogram:methods]{see table} \\
regular       & \hyperref[regular:methods]{see table}       \\
vector        & \hyperref[vector:methods]{see table}        \\
matrix        & \hyperref[matrix:methods]{see table}        \\
angle         & \hyperref[angle:methods]{see table}         \\
list\_point   & \hyperref[list_point:methods]{see table}    \\
fct           & \hyperref[fct:methods]{see table}           \\
pfct          & \hyperref[pfct:methods]{see table}          \\
\end{tabular}
\section*{Metamethods Overview}
\begin{tabular}{ll}
\textbf{Class} & \textbf{Methods} \\
\hline
point  & \hyperref[point:methods]{see table}  \\
occs   & \hyperref[occs:methods]{see table}   \\
vector & \hyperref[vector:methods]{see table} \\
matrix & \hyperref[matrix:methods]{see table} \\
path   & \hyperref[path:methods]{see table}   \\
\end{tabular}
\clearpage
\tableofcontents

\clearpage
\newpage

\part{Geometry Core}

\section{Getting started}

\subsection{The first code}

A quick introduction to get you started.
We assume that the packages \tkzNamePack{tkz-euclide} and \tkzNamePack{tkz-elements} are installed.
Compile the following code using the \tkzEngine{lualatex} engine; you should obtain a straight line
passing through points $A$ and $B$.

\medskip
\noindent
\textbf{Note.}
The package \tkzNamePack{tkz-elements} performs all geometric definitions and computations in Lua,
while \tkzNamePack{tkz-euclide} is mainly used for drawing.
For this reason, the \texttt{mini} option of \tkzNamePack{tkz-euclide} is recommended.
If a compilation problem occurs, simply load \tkzNamePack{tkz-euclide} without this option.


\begin{mybox}{}
\begin{minipage}{.55\textwidth}
  \begin{tkzexample}[code only]
  % !TEX TS-program = lualatex
  \documentclass{article}
  \usepackage[mini]{tkz-euclide}
  \usepackage{tkz-elements}
  \begin{document}
  \directlua{
    init_elements()
    z.A = point(0, 1)
    % or z.A = point:new(0, 1)
    z.B = point(2, 0)
  }
  \begin{tikzpicture}
   \tkzGetNodes
   \tkzDrawLine(A,B)
   \tkzDrawPoints(A,B)
   \tkzLabelPoints(A,B)
  \end{tikzpicture}
  \end{document}
  \end{tkzexample}
\end{minipage}
\begin{minipage}{.35\textwidth}
  \directlua{
    z.A = point(0, 1)
    z.B = point(2, 0)
  }
  \begin{tikzpicture}
   \tkzGetNodes
   \tkzDrawLine(A,B)
   \tkzDrawPoints(A,B)
   \tkzLabelPoints(A,B)
  \end{tikzpicture}
\end{minipage}
\end{mybox}

\subsection{Key points}

The following points are essential for most of the codes in this documentation.

\begin{itemize}
\item \verb|% !TEX TS-program = lualatex|.
Compilation must be performed with \code{lualatex}. This line is optional if your editor is already configured accordingly.

\item \verb|\usepackage[mini]{tkz-euclide}|.
The use of \tkzNamePack{tkz-euclide} is optional; however, if it is loaded, the \code{mini} option is recommended.
It loads only the drawing macros, while all computations are handled by \code{lua}.
\textbf{Important:} at the current stage, some drawing macros are not yet fully independent of computation macros.
If a compilation problem occurs, simply load \tkzNamePack{tkz-euclide} without the \code{mini} option.

\item \verb|\usepackage{tkz-elements}|.
This package is required. It provides the Lua-based geometry engine and support macros used together with \tkzNamePack{tkz-euclide}.

\item \verb|\directlua{...}|.
All geometric definitions and computations are placed inside this macro (or within a \texttt{tkzelements} environment).

\item \verb|init_elements()|.
This function should be called at the beginning of each Lua section.
It resets internal tables and clears global variables used by \tkzNamePack{tkz-elements}.

\item \verb|\tkzGetNodes|.
This macro must be placed at the beginning of the \code{tikzpicture} environment.
It transfers the points defined in Lua to TikZ as usable nodes.
\end{itemize}



\subsection{Testing}

To test your installation and follow the examples in this documentation, you need to load two packages: \tkzNamePack{tkz-euclide} and \tkzNamePack{tkz-elements}. The first package automatically loads \tkzNamePack{\TIKZ}, which is necessary for all graphical rendering.


The \tkzname{Lua} code is provided as an argument to the \tkzMacro{lualatex}{directlua} macro. We will often refer to this code block as the \tkzname{Lua part } \footnote{This code can also be placed in an external file, e.g., \texttt{file.lua}.}. This part depends entirely on the \tkzNamePack{tkz-elements} package.

A crucial component in the \tkzEnv{tikz}{tikzpicture} environment is the macro \tkzMacro{tkz-elements}{tkzGetNodes}. This macro transfers the points defined in \tkzname{Lua} to \tkzNamePack{\TIKZ} by creating the corresponding nodes. All such points are stored in a table named \tkzname{z} and are accessed using the syntax \tkzname{z.label}. These labels are then reused within \tkzNamePack{tkz-euclide}.

When you define a point by assigning it a label and coordinates, it is internally represented as a complex number — the affix of the point. This representation allows the point to be located within an orthonormal Cartesian coordinate system.

If you want to use a different method for rendering your objects, this is the macro to modify. For example, Section~\ref{sec:metapost} presents \tkzMacro{tkz-elements}{tkzGetNodesMP}, a variant that enables communication with \code{MetaPost}.

Another essential element is the use of the function \tkzFct{tkz-elements}{init\_elements()}, which clears internal tables\footnote{All geometric objects are stored in Lua tables. These tables must be cleaned regularly, especially when creating multiple figures in sequence.} when working with multiple figures.

If everything worked correctly with the previous code, you're ready to begin creating geometric objects. Section~\ref{sec:class_and_object} introduces the available options and object structures.

Finally, it is important to be familiar with basic drawing commands in \tkzNamePack{tkz-euclide}, as they will be used to render the objects defined in \tkzname{Lua}.
\endinput
\input{TKZdoc-elements-structure.tex}
\input{TKZdoc-elements-why.tex}
\input{TKZdoc-elements-presentation.tex}
\input{TKZdoc-elements-convention.tex}
\input{TKZdoc-elements-organization.tex}
\input{TKZdoc-elements-coordinates.tex}
%--------------------------------------------------------------------
\section{Numerical Tolerance}
\label{sec:numerical_tolerance}

\subsection{Floating-Point Arithmetic}

All computations in \tkzNamePack{tkz-elements} are performed using
floating-point arithmetic.
As a consequence, exact comparisons between real numbers are unreliable.

For example, a point theoretically lying on a line may produce a
very small non-zero value due to rounding errors.

\subsection{Global Tolerance: \code{tkz.epsilon}}

To ensure numerical robustness, \tkzNamePack{tkz-elements}
uses a global tolerance parameter:


\verb|tkz.epsilon = 1e-10|


This value defines the admissible numerical error in geometric tests
(collinearity, incidence, equality of distances, etc.).

\medskip
By default, this tolerance is set to a small positive value.
It can be adjusted by advanced users if needed.

\subsection{Usage in Position Tests}

All membership and position tests are \emph{EPS-aware}.
When a method accepts an optional argument \code{EPS}, the following rule applies:

\begin{itemize}
\item If \code{EPS} is provided, it overrides the global tolerance.
\item Otherwise, \code{tkz.epsilon} is used.
\end{itemize}

This design ensures consistency across all geometric objects.
%--------------------------------------------------------------------
\section{Geometric Relations API}
\label{sec:geometric_relations_api}

\subsection{General Principle}

In \tkzNamePack{tkz-elements}, geometric objects provide a unified method
\[
\code{object:position(other\_object[, EPS])}
\]
to classify a geometric relation.

The result depends on the type of the second argument (\code{other\_object}).

\medskip
All position tests are tolerance-aware.
If not specified, the optional parameter \code{EPS} defaults to the global value
\tkzVar{tkz}{tkz.epsilon}.

\medskip
\textbf{Convention:}
All classification results are returned as \emph{uppercase symbolic strings},
intended for logical comparison in Lua code.

%--------------------------------------------------------------------
\subsection{Point vs Object}

When the tested object is a point, the result describes membership relative to
a region (or a boundary).

\medskip
Possible return values:

\begin{center}
\begin{tabular}{ll}
\code{"IN"}  & strictly inside the region, \\
\code{"ON"}  & on the boundary (within tolerance), \\
\code{"OUT"} & strictly outside the region.
\end{tabular}
\end{center}

\medskip
Examples:
\begin{itemize}
\item \code{circle:position(point)} \quad (\code{"IN"}, \code{"ON"}, \code{"OUT"})
\item \code{triangle:position(point)} \quad (\code{"IN"}, \code{"ON"}, \code{"OUT"})
\item \code{conic:position(point)} \quad (\code{"IN"}, \code{"ON"}, \code{"OUT"})
\item \code{line:position(point)} \quad (\code{"ON"} or \code{"OUT"})
\end{itemize}

%--------------------------------------------------------------------
\subsection{Object vs Object}

When both arguments are geometric objects, the result describes their mutual
geometric relation.

\subsubsection*{Line vs Line}

\begin{center}
\begin{tabular}{ll}
\code{"INTERSECT"} & lines intersect at one point, \\
\code{"PARALLEL"}  & distinct parallel lines, \\
\code{"IDENTICAL"} & identical lines.
\end{tabular}
\end{center}

\subsubsection*{Line vs Circle}

\begin{center}
\begin{tabular}{ll}
\code{"DISJOINT"} & no intersection, \\
\code{"TANGENT"}  & exactly one common point, \\
\code{"SECANT"}   & two intersection points.
\end{tabular}
\end{center}

\subsubsection*{Circle vs Circle}

\begin{center}
\begin{tabular}{ll}
\code{"DISJOINT\_EXT"} & exterior disjoint circles, \\
\code{"TANGENT\_EXT"}  & exterior tangency, \\
\code{"SECANT"}        & two intersection points, \\
\code{"TANGENT\_INT"}  & interior tangency, \\
\code{"DISJOINT\_INT"} & one circle strictly inside the other, \\
\code{"CONCENTRIC"}    & same center, different radii, \\
\code{"IDENTICAL"}     & identical circles.
\end{tabular}
\end{center}

%--------------------------------------------------------------------
\subsection{Triangle and Conic}

For \tkzClass{triangle} and \tkzClass{conic}, the method \code{position()}
currently supports point arguments only.

\subsubsection*{Triangle vs Point}

\begin{center}
\begin{tabular}{ll}
\code{"IN"}  & strictly inside the triangle, \\
\code{"ON"}  & on an edge or a vertex, \\
\code{"OUT"} & outside the triangle.
\end{tabular}
\end{center}

\subsubsection*{Conic vs Point}

\begin{center}
\begin{tabular}{ll}
\code{"IN"}  & inside the conic region (according to its type), \\
\code{"ON"}  & on the conic curve, \\
\code{"OUT"} & outside the associated region.
\end{tabular}
\end{center}

\medskip
\texttt{Note:} For hyperbolas and parabolas, the meaning of \code{"IN"} and
\code{"OUT"} depends on the region associated with the conic type and its
construction data (focus/directrix or equivalent representation).

%--------------------------------------------------------------------
\subsection{Design Philosophy}

\begin{itemize}
\item A single polymorphic method centralizes geometric classification.
\item Uppercase symbolic results ensure clarity and stability.
\item Numerical tolerance improves robustness near boundary cases.
\item The API is extensible: additional object types may be supported over time.
\item Backward compatibility wrappers may be provided when older boolean methods
      existed in previous versions.
\end{itemize}
\input{TKZdoc-elements-classes.tex}
\input{TKZdoc-elements-classes-point.tex}
\input{TKZdoc-elements-classes-line.tex}
\input{TKZdoc-elements-classes-circle.tex}
\input{TKZdoc-elements-classes-triangle.tex}
\input{TKZdoc-elements-classes-occs.tex}
\input{TKZdoc-elements-classes-conic.tex}
\newpage
\section{Class \tkzClass{quadrilateral}}

\vspace{1em}
The variable \tkzVar{quadrilateral}{Q} holds a table used to store quadrilaterals. It is optional, and you are free to choose the variable name. However, using \code{Q} is a recommended convention for clarity and consistency. If you use a custom variable (e.g., Quad), you must initialize it manually. The \tkzFct{tkz-elements}{init\_elements()} function reinitializes the \code{Q} table if used.


\subsection{Creating a quadrilateral}
\label{sub:creating_a_quadrilateral}
The \tkzClass{quadrilateral} class requires four points. The order defines the sides of the quadrilateral.

\medskip
The object is usually stored in \tkzVar{quadrilateral}{Q}, which is the recommended variable name.

\begin{mybox}
\begin{verbatim}
Q.ABCD = quadrilateral:new(z.A, z.B, z.C, z.D)
\end{verbatim}
\end{mybox}

\textbf{Short form.}

A shorter syntax is also available:

\begin{mybox}
\begin{verbatim}
Q.ABCD = quadrilateral(z.A, z.B, z.C, z.D)
\end{verbatim}
\end{mybox}

\subsection{Quadrilateral Attributes}
Points are created in the direct direction. A test is performed to check whether the points form a rectangle, otherwise compilation is blocked.

\begin{center}
 \code{ Creation : Q.new = rectangle : new (z.A,z.B,z.C,z.D)}
\end{center}


  \bgroup
  \catcode`_=12
  \small
  \captionof{table}{rectangle attributes.}\label{quadrilateral:attributes}
  \begin{tabular}{lll}
  \toprule
  \textbf{Attributes}         & \textbf{Application} & \\
  \tkzAttr{quadrilateral}{pa}   & |z.A = Q.new.pa| & \\
  \tkzAttr{quadrilateral}{pb}   & |z.B = Q.new.pb| & \\
  \tkzAttr{quadrilateral}{pc}   & |z.C = Q.new.pc| & \\
  \tkzAttr{quadrilateral}{pd}   & |z.D = Q.new.pd| & \\
  \tkzAttr{quadrilateral}{type} & |Q.new.type= 'quadrilateral'|  &\\
  \tkzAttr{quadrilateral}{center}    & |z.I = Q.new.center| & intersection of diagonals\\
  \tkzAttr{quadrilateral}{g}    & |z.G = Q.new.g| & barycenter\\
  \tkzAttr{quadrilateral}{a}    & |AB = Q.new.a| & barycenter\\
  \tkzAttr{quadrilateral}{b}    & |BC = Q.new.b| & barycenter\\
  \tkzAttr{quadrilateral}{c}    & |CD = Q.new.c| & barycenter\\
  \tkzAttr{quadrilateral}{d}    & |DA = Q.new.d| & barycenter\\
  \tkzAttr{quadrilateral}{ab}   & |Q.new.ab|   &  line passing through two vertices   \\
  \tkzAttr{quadrilateral}{ac}   & |Q.new.ca|   &  idem. \\
  \tkzAttr{quadrilateral}{ad}   & |Q.new.ad|   &  idem. \\
  \tkzAttr{quadrilateral}{bc}   & |Q.new.bc|   &  idem. \\
  \tkzAttr{quadrilateral}{bd}   & |Q.new.bd|   &  idem. \\
  \tkzAttr{quadrilateral}{cd}   & |Q.new.cd|   &  idem. \\
  \bottomrule
  \end{tabular}
  \egroup


\subsubsection{Quadrilateral attributes}

\vspace{1em}
\begin{minipage}{.5\textwidth}
\begin{verbatim}
\directlua{
init_elements()
z.A = point(0, 0)
z.B = point(4, 0)
z.C = point(5, 1)
z.D = point(0, 3)
Q.ABCD = quadrilateral(z.A, z.B, z.C, z.D)
z.I = Q.ABCD.i
z.G = Q.ABCD.g}

\begin{tikzpicture}
\tkzGetNodes
\tkzDrawPolygon(A,B,C,D)
\tkzDrawSegments(A,C B,D)
\tkzDrawPoints(A,B,C,D,I,G)
\end{tikzpicture}
\end{verbatim}
\end{minipage}
\begin{minipage}{.5\textwidth}
\directlua{
 init_elements()
 z.A = point(0, 0)
 z.B = point(4, 0)
 z.C = point(5, 1)
 z.D = point(0, 3)
 Q.ABCD = quadrilateral(z.A, z.B, z.C, z.D)
 z.I = Q.ABCD.center
 z.G = Q.ABCD.g}

\begin{center}
  \begin{tikzpicture}
  \tkzGetNodes
  \tkzDrawPolygon(A,B,C,D)
  \tkzDrawSegments(A,C B,D)
  \tkzDrawPoints(A,B,C,D,I,G)
  \end{tikzpicture}
\end{center}
\end{minipage}

\subsubsection{Quadrilateral examples}
Advanced Euclidean Geometry 2013 Supplement June 26\\

\texttt{Construction of geometric mean: }


\directlua{
init_elements()
z.A = point(0, 0)
z.B = point(1, 4)
z.C = point(3.5, 4)
z.D = point(5, 0 )
Q.ABCD = quadrilateral(z.A,z.B,z.C,z.D)
z.I = Q.ABCD.center
L.ll = Q.ABCD.da:ll_from (z.I)
z.H = intersection(L.ll,Q.ABCD.ab)
z.Hp = intersection(L.ll,Q.ABCD.cd)
z.M =  Q.ABCD.ab.mid
C.MA = circle(z.M,z.A)
L.perp = Q.ABCD.ab:ortho_from(z.H)
_,z.P = intersection (C.MA,L.perp)
T.PAB = triangle (z.P,z.A,z.B)
L.bis =  T.PAB:bisector ()
z.G = L.bis.pb
L.llg = Q.ABCD.da: ll_from (z.G)
L.llm = Q.ABCD.da : ll_from (z.M)
z.Gp = intersection (L.llg,Q.ABCD.cd)}
\begin{center}
\begin{tikzpicture}
\tkzGetNodes
\tkzDrawPolygon(A,B,C,D)
\tkzDrawCircle(M,A)
\tkzDrawSegments[red](H,H')
\tkzDrawSegments[blue](G,G')
\tkzDrawSegments[dashed](P,H P,B P,A P,G A,C B,D)
\tkzDrawPoints(A,B,C,D,H,I,M,P,G,G',H')
\tkzLabelPoints(A,D,I)
\tkzLabelPoints[above](B,C,P)
\tkzLabelPoints[above right=-3pt](H,G)
\tkzLabelPoints[right](G',H')
 \tkzLabelSegment[above](B,C){$b$}
 \tkzLabelSegment[below](A,D){$a$}
\end{tikzpicture}
\end{center}



\begin{proof}
A trapezoid with parallel sides $AD$ and $BC$ is given.

Let H be the endpoint of the parallel through $I$
to the side $AD$.  $HH' = \dfrac{2ab}{a+b}$ .

$PH$ is perpendicular to $AB$

$(PG)$ is the bisector of $\widehat{APB}$

 $GG' = \dfrac{a \cdot \sqrt{b} + b \cdot \sqrt{a}}{\sqrt{a}+\sqrt{b}}
= \dfrac{\sqrt{ab}(\sqrt{a}+\sqrt{b})}{\sqrt{a}+\sqrt{b}} = \sqrt{ab} $
\end{proof}

\texttt{Code: }

\begin{verbatim}
\directlua{
init_elements()
z.A = point(0, 0)
z.B = point(1, 4)
z.C = point(3.5, 4)
z.D = point(5, 0 )
Q.ABCD = quadrilateral(z.A,z.B,z.C,z.D)
z.I = Q.ABCD.center
L.ll = Q.ABCD.da:ll_from (z.I)
z.H = intersection(L.ll,Q.ABCD.ab)
z.Hp = intersection(L.ll,Q.ABCD.cd)
z.M =  Q.ABCD.ab.mid
C.MA = circle(z.M,z.A)
L.perp = Q.ABCD.ab:ortho_from(z.H)
_,z.P = intersection (C.MA,L.perp)
T.PAB = triangle (z.P,z.A,z.B)
L.bis =  T.PAB:bisector ()
z.G = L.bis.pb
L.llg = Q.ABCD.da: ll_from (z.G)
L.llm = Q.ABCD.da : ll_from (z.M)
z.Gp = intersection (L.llg,Q.ABCD.cd)}
\end{verbatim}

\subsection{Quadrilateral methods}

\vspace{1em}

  \bgroup
  \catcode`_=12
  \small
  \captionof{table}{Quadrilateral methods.}\label{quadrilateral:methods}
  \begin{tabular}{ll}
  \toprule
  \textbf{Methods} & \textbf{Reference}    \\
  \midrule
  \tkzMeth{quadrilateral}{new()}& Note\footnote{quadrilateral(pt, pt, pt, pt) (short form, recommended)}; [\ref{sub:creating_a_quadrilateral}]\\

  \tkzMeth{quadrilateral}{is\_cyclic ()} & [\ref{ssub:iscyclic}]\\
  \tkzMeth{quadrilateral}{is\_convex ()} & [\ref{ssub:isconvex}]\\
  \tkzMeth{quadrilateral}{poncelet\_point()} & [\ref{ssub:poncelet_point}]\\
  \bottomrule %
  \end{tabular}
  \egroup



\subsubsection{Method \tkzMeth{quadrilateral}{is\_cyclic()}}
\label{ssub:iscyclic}

\textbf{Inscribed quadrilateral}

\vspace{1em}
\begin{minipage}{.5\textwidth}
\begin{verbatim}
\directlua{
 init_elements()
 z.A = point(0, 0)
 z.B = point(4, 0)
 z.D = point:polar(4, tkz.tau / 3)
 L.DB = line(z.D, z.B)
 T.equ = L.DB:equilateral()
 z.C = T.equ.pc
 Q.new = quadrilateral(z.A, z.B, z.C, z.D)
 bool = Q.new:is_cyclic()
 if bool == true then
 	C.cir = triangle(z.A, z.B, z.C):
   circum_circle()
 	z.O = C.cir.center
 end}
\begin{tikzpicture}
  \tkzGetNodes
  \tkzDrawPolygon(A,B,C,D)
  \tkzDrawPoints(A,B,C,D)
  \tkzDrawCircle(O,A)
  \ifthenelse{\equal{\tkzUseLua{bool}}{
     true}}{\tkzDrawCircle(O,A)}{}
  \tkzLabelPoints(A,B)
  \tkzLabelPoints[above](C,D)
\end{tikzpicture}
\end{verbatim}
\end{minipage}
\begin{minipage}{.5\textwidth}
\directlua{
init_elements()
z.A = point(0, 0)
z.B = point(4, 0)
z.D = point:polar(4, tkz.tau / 3)
L.DB = line(z.D, z.B)
T.equ = L.DB:equilateral()
z.C = T.equ.pc
Q.new = quadrilateral(z.A, z.B, z.C, z.D)
bool = Q.new:is_cyclic()
if bool == true then
	C.cir = triangle(z.A, z.B, z.C):circum_circle()
	z.O = C.cir.center
end
}
\begin{center}
\begin{tikzpicture}[scale = .75 ]
  \tkzGetNodes
  \tkzDrawPolygon(A,B,C,D)
  \tkzDrawPoints(A,B,C,D)
  \tkzDrawCircle(O,A)
  \ifthenelse{\equal{\tkzUseLua{bool}}{
   true}}{\tkzDrawCircle(O,A)}{}
  \tkzLabelPoints(A,B)
  \tkzLabelPoints[above](C,D)
\end{tikzpicture}
\end{center}
\end{minipage}

\subsubsection{Method \tkzMeth{quadrilateral}{is\_convex()}}
\label{ssub:isconvex}







\subsubsection{Method \tkzMeth{quadrilateral}{poncelet\_point}}
\label{ssub:poncelet_point}

See [\ref{ssub:method_triangle_poncelet_point}] for the definition.

\begin{tkzexample}[latex=.5\textwidth]
\directlua{
 init_elements()
 z.A = point(1, 1)
 z.B = point(6, 0)
 z.D = point(0, 5)
 z.C = point(6, 2)
 Q.ABCD = quadrilateral(z.A, z.B, z.C, z.D)
 z.P = Q.ABCD:poncelet_point()
 T.ABC = triangle(z.A,z.B,z.C)
 z.I = T.ABC.eulercenter
 z.Mc = tkz.midpoint(z.A, z.B)
 T.ABD = triangle(z.A, z.B, z.D)
 z.I1 = T.ABD.eulercenter}
\begin{center}
\begin{tikzpicture}
  \tkzGetNodes
  \tkzDrawPolygons(A,B,C,D)
  \tkzDrawPoints(A,B,C,D,P,Mc)
  \tkzDrawCircles[red](I,Mc I1,Mc)
  \tkzLabelPoints(A,B)
  \tkzLabelPoints[above](C,D,P)
\end{tikzpicture}
\end{center}
\end{tkzexample}
\endinput
\input{TKZdoc-elements-classes-square.tex}
\input{TKZdoc-elements-classes-rectangle.tex}
\input{TKZdoc-elements-classes-parallelogram.tex}
\input{TKZdoc-elements-classes-regular.tex}

\part{Algebra and Tools}

\input{TKZdoc-elements-classes-vectors.tex}
\input{TKZdoc-elements-classes-matrices.tex}
 \input{TKZdoc-elements-classes-path.tex}
\section{Class \tkzClass{list\_point}}
\label{sec:class_list_point}

The variable \tkzVar{list\_point}{LP} holds a table used to store \tkzClass{list\_point} objects.
Its use is optional, and users are free to choose any variable name.

\medskip

However, for the sake of consistency and readability throughout the documentation,
it is recommended to use the predefined variable \tkzVar{list\_point}{LP}.
The function \tkzFct{tkz-elements}{init\_elements()} automatically reinitializes this table.

\medskip

The \tkzClass{list\_point} class represents an ordered list of points.
It is primarily intended as a lightweight container used to store and manipulate multiple points
produced by geometric constructions, intersections, or iterative algorithms.

Unlike geometric objects such as lines or circles, a \tkzClass{list\_point} object does not define intrinsic geometry.
Its purpose is to facilitate grouping, iteration, and data transfer between \tkzname{Lua} and \tkzname{TikZ}.

\medskip

A \tkzClass{list\_point} object behaves like a Lua table indexed by integers,
while providing a set of convenience methods adapted to geometric workflows.

\medskip

\subsection{Attributes}
\label{subsec:list_point_attributes}

\vspace{1em}
\bgroup
\small
\captionof{table}{\texttt{list\_point} attributes.}\label{list_point:attributes}
\begin{tabular}{lll}
\toprule
\textbf{Attribute} & \textbf{Meaning} & \texttt{Reference} \\
\midrule
\tkzAttr{list\_point}{n} &  [\ref{ssub:attr_list_point_n}] \\
\tkzAttr{list\_point}{items} & [\ref{ssub:attr_list_point_items}] \\
\bottomrule
\end{tabular}
\egroup


\subsubsection{Attribute \tkzAttr{list\_point}{n}}
\label{ssub:attr_list_point_n}

The attribute \tkzAttr{list\_point}{n} denotes the number of points stored in the list.
In practice, it corresponds to the Lua length operator \texttt{\#P} (array part).

\subsubsection{Attribute \tkzAttr{list\_point}{items}}
\label{ssub:attr_list_point_items}

The attribute \tkzAttr{list\_point}{items} refers to the ordered sequence of points stored in the list.
Individual points are accessed using standard Lua indexing, e.g. \texttt{P[i]}.


\subsection{Creating an object}
\label{subsec:list_point_create}

\subsubsection{Method \tkzMeth{list\_point}{new(...)}}
\label{ssub:meth_list_point_new}

Creates a new \tkzClass{list\_point} object.

In user code, the recommended constructor is the callable form \code{list\_point(...)}.
Internally, this calls the method \tkzMeth{list\_point}{new(...)}.

\medskip

A \tkzClass{list\_point} object can be created either as a local Lua variable
or as an entry of the table \tkzVar{list\_point}{LP}, which is reserved for storing
lists of points.

\begin{tkzexample}[code only]
local P = list_point()
LP.Q = list_point(z.A, z.B, z.C)
\end{tkzexample}

Both forms are valid and produce independent \tkzClass{list\_point} objects.

\medskip

Using a local variable is convenient for temporary computations or intermediate results.
However, for objects intended to be reused, transferred to \tkzname{TikZ},
or accessed outside a local scope, it is recommended to store them in the
\tkzVar{list\_point}{LP} table.

\medskip

The function \tkzFct{tkz-elements}{init\_elements()} automatically reinitializes
the \tkzVar{list\_point}{LP} table.


% --------------------------------------------------------------------
\subsection{Methods or Basic accessors}
\label{subsec:list_point_accessors}

\vspace{1em}
\bgroup
\small
\captionof{table}{\texttt{list\_point} methods.}\label{list_point:methods}
\begin{tabular}{ll}
\toprule
\texttt{Method}  & \texttt{Reference} \\
\midrule
\texttt{Constructor} &    \\
\midrule
\tkzMeth{list\_point}{new(...)}   & [\ref{ssub:meth_list_point_new}] \\
\midrule
\texttt{Methods Returning a Integer Number}&\\
\midrule
\tkzMeth{list\_point}{len()}      & [\ref{ssub:meth_list_point_len}] \\
\midrule
\texttt{Methods Returning a Point}&\\
\midrule\tkzMeth{list\_point}{get(i)}     & [\ref{ssub:meth_list_point_get}] \\
\tkzMeth{list\_point}{barycenter()}& [\ref{ssub:meth_list_point_barycenter}]\\
\midrule
\texttt{Methods Returning a List\_Point}&\\
\midrule
\tkzMeth{list\_point}{add(p)}     & [\ref{ssub:meth_list_point_add}] \\
\tkzMeth{list\_point}{extend(pl)} & [\ref{ssub:meth_list_point_extend}] \\
\tkzMeth{list\_point}{map(f)}     & [\ref{ssub:meth_list_point_map}] \\
\midrule
\tkzMeth{list\_point}{unpack()}   & [\ref{ssub:meth_list_point_unpack}] \\


\tkzMeth{list\_point}{clear()}    & [\ref{ssub:meth_list_point_clear}] \\
\tkzMeth{list\_point}{foreach(f)} & [\ref{ssub:meth_list_point_foreach}] \\

\tkzMeth{list\_point}{bbox()}     & [\ref{ssub:meth_list_point_bbox}] \\

\tkzMeth{list\_point}{to\_path()}& [\ref{ssub:meth_list_point_to_path}] \\
\bottomrule
\end{tabular}
\egroup

\subsubsection{Method \tkzMeth{list\_point}{len()}}
\label{ssub:meth_list_point_len}

Returns the number of points stored in the list.

\subsubsection{Method \tkzMeth{list\_point}{get(i)}}
\label{ssub:meth_list_point_get}

Returns the point at index \code{i}.

\subsubsection{Method \tkzMeth{list\_point}{unpack()}}
\label{ssub:meth_list_point_unpack}

Returns all points as separate values using \code{table.unpack}.
This method is useful when interfacing with functions expecting individual point arguments.

\begin{tkzexample}[latex=.5\textwidth]
\directlua{
init_elements()
LP = list_point()
LP:add(point(0,0))
LP:add(point(1,0))
LP:add(point(1,1))
z.A, z.B, z.C = LP:unpack()
}
\begin{center}
\begin{tikzpicture}
\tkzGetNodes
\tkzDrawPolygon(A,B,C)
\end{tikzpicture}
\end{center}
\end{tkzexample}

% --------------------------------------------------------------------
\subsection{Modification methods}
\label{subsec:list_point_mutators}

\subsubsection{Method \tkzMeth{list\_point}{add(p)}}
\label{ssub:meth_list_point_add}

Appends the point \code{p} to the list and returns the object itself.

\subsubsection{Method \tkzMeth{list\_point}{extend(pl)}}
\label{ssub:meth_list_point_extend}

Appends all points from another \tkzClass{list\_point} object \code{pl}.

\subsubsection{Method \tkzMeth{list\_point}{clear()}}
\label{ssub:meth_list_point_clear}

Removes all points from the list.


These methods modify the object in place and allow method chaining.

% --------------------------------------------------------------------
\subsection{Iteration helpers}
\label{subsec:list_point_iteration}

\subsubsection{Method \tkzMeth{list\_point}{foreach(f)}}
\label{ssub:meth_list_point_foreach}

Applies the function \code{f(p, i)} to each point \code{p} with its index \code{i}.

\subsubsection{Method \tkzMeth{list\_point}{map(f)}}
\label{ssub:meth_list_point_map}

Applies the function \code{f(p, i)} to each point and returns a new \tkzClass{list\_point} object
containing the results.


These methods support a functional programming style for geometric transformations.

% --------------------------------------------------------------------
\subsection{Geometric utilities}
\label{subsec:list_point_geometry}

\subsubsection{Method \tkzMeth{list\_point}{barycenter()}}
\label{ssub:meth_list_point_barycenter}

Returns the barycenter (centroid) of all points in the list.
If the list is empty, the method returns \code{nil}.

\subsubsection{Method \tkzMeth{list\_point}{bbox()}}
\label{ssub:meth_list_point_bbox}

Returns the bounding box of the point set as four numbers:
\[
  (x_{\min}, y_{\min}, x_{\max}, y_{\max})
\]


% --------------------------------------------------------------------
\subsection{TikZ output}
\label{subsec:list_point_output}

\subsubsection{Method \tkzMeth{list\_point}{to\_path()}}
\label{ssub:meth_list_point_to_path}

Converts the list of points into a TikZ path string of the form
\[
  (x_1,y_1) -- (x_2,y_2) -- \dots
\]
suitable for drawing with TikZ commands.


\medskip

This class provides a minimal yet effective bridge between \tkzname{Lua} computations and \tkzname{TikZ} drawing instructions.


\subsection{Examples}
\label{subsec:list_point_examples}

% --------------------------------------------------------------------
\subsubsection{Temporary list of points}
\label{ssub:list_point_example_local}

This example illustrates the use of a local \tkzClass{list\_point} object to store
temporary results inside a computation.

\begin{tkzexample}[code only]
local P = list_point()
for i = 1, 10 do
  P:add(z.A + i * (z.B - z.A) / 10)
end
local G = P:barycenter()
\end{tkzexample}

The list \code{P} is used only locally to compute the barycenter of a set of points.

% --------------------------------------------------------------------
\subsubsection{Storing construction results in \tkzVar{list\_point}{LP}}
\label{ssub:list_point_example_lp}

This example shows how to store a list of intersection points for later use.

\begin{tkzexample}[code only]
LP.I = list_point()
local X, Y = intersection(L.AB, C.circ)
LP.I:add(X)
LP.I:add(Y)
\end{tkzexample}

Storing the points in \tkzVar{list\_point}{LP} makes them accessible for drawing
or further constructions.

% --------------------------------------------------------------------
\subsubsection{Iterative construction}
\label{ssub:list_point_example_iterative}

A typical use case is the progressive construction of a family of points.

\begin{tkzexample}[code only]
LP.P = list_point()
for i = 0, 24 do
  local t = i / 24
  LP.P:add(L.AB:point(t))
end
\end{tkzexample}

The list \code{LP.P} contains a discrete sampling of a segment.

% --------------------------------------------------------------------
\subsubsection{Functional transformation}
\label{ssub:list_point_example_map}

This example uses the \tkzMeth{list\_point}{map} method to transform a set of points.

\begin{tkzexample}[code only]
local P = list_point(z.A, z.B, z.C)
LP.Q = P:map(function(p)
  return rotation_(p, tkz.tau / 3, z.O)
end)
\end{tkzexample}

The original list \code{P} is preserved, while \code{LP.Q} stores the transformed points.

% --------------------------------------------------------------------
\subsubsection{Conversion to a TikZ path}
\label{ssub:list_point_example_path}

This example illustrates the interaction between \tkzClass{list\_point} and TikZ.

\begin{tkzexample}[code only]
LP.P = list_point(z.A, z.B, z.C)
local path = LP.P:to_path()
\end{tkzexample}

The string returned by \tkzMeth{list\_point}{to\_path} can be passed directly
to TikZ drawing commands.


% --------------------------------------------------------------------
\subsubsection{From \tkzClass{list\_point} to \tkzClass{path} and drawing with \tkzNamePack{tkz-euclide}}
\label{ssub:list_point_example_points_path_draw}

This example shows a typical workflow: compute a list of points in \tkzname{Lua},
convert it into a \tkzClass{path}, and then draw the resulting points with
\tkzMeth{tkz-euclide}{tkzDrawPointsFromPath}.


\begin{tkzexample}[latex = .5\textwidth]
\directlua{
  init_elements()
  z.A = point(0, 0)
  z.B = point(8, 6)
  L.AB = line(z.A, z.B)
  % list of sampled points
  LP.P = list_point()
  for i = 0, 12 do
    LP.P:add(L.AB:point(i/12))
  end
  % convert list_point -> path
  PA.sample = path()
  for i = 1, LP.P:len() do
    PA.sample:add_point(LP.P:get(i))
  end
}
\begin{center}
\begin{tikzpicture}
\tkzGetNodes
% draw the points stored in the path
\tkzDrawPointsFromPath(PA.sample)
\end{tikzpicture}
\end{center}
\end{tkzexample}

The points are stored in \code{LP.P} for later reuse, while \code{PA.sample}
contains the corresponding TikZ path representation used for drawing.

\newpage
\section{Class \tkzClass{angle}}
\label{sec:class_angle}

The \tkzClass{angle} class is an experimental helper object used to represent
an angle defined by three points.
It is currently self-contained and does not interact with other classes.
Its main purpose is to provide a simple and direct interface for obtaining:

\begin{itemize}
  \item the oriented angle in radians,
  \item the normalised oriented angle in $[0, 2\pi]$
  \item the interior (non-oriented) angle,
  \item the measure of the interior angle in degrees.
\end{itemize}

An \tkzClass{angle} object is \emph{static}: all values are computed at creation
time and never updated.

\subsection{Creating an object}
\label{ssub:method_angle_new}

\begin{verbatim}
local angle = require("tkz_elements_angle")

local alpha  = angle(A, B, C)   -- equivalent to angle:new(A,B,C)
local beta   = angle(B, C, A)
local gamma  = angle(C, A, B)
\end{verbatim}

The three arguments are:
\begin{itemize}
  \item \tkzVar{ps} : the vertex of the angle,
  \item \tkzVar{pa} : the first point defining the first ray,
  \item \tkzVar{pb} : the second point defining the second ray.
\end{itemize}

\subsection{Attributes}

The following attributes are stored inside every \tkzClass{angle} object:

\vspace{1em}
\bgroup
\small
\captionof{table}{Angle attributes.}\label{angle:attributes}
\begin{tabular}{lll}
\toprule
\textbf{Attribute} & \textbf{Meaning} & \textbf{Reference} \\
\midrule
\tkzAttr{angle}{ps}   & Vertex of the angle & [\ref{ssub:attr_angle_ps}] \\
\tkzAttr{angle}{pa}   & First defining point (ray $[ps\,pa]$) & [\ref{ssub:attr_angle_pa}] \\
\tkzAttr{angle}{pb}   & Second defining point (ray $[ps\,pb]$) & [\ref{ssub:attr_angle_pb}] \\
\tkzAttr{angle}{raw}  & Oriented angle (radians), may be negative & [\ref{ssub:attr_angle_raw}] \\
\tkzAttr{angle}{norm} & Oriented angle normalised to $[0,2\pi)$ & [\ref{ssub:attr_angle_norm}] \\
\bottomrule
\end{tabular}
\egroup

\paragraph{\tkzAttr{angle}{ps}}\label{ssub:attr_angle_ps}
The vertex (summit) of the angle.

\paragraph{\tkzAttr{angle}{pa}}\label{ssub:attr_angle_pa}
A point defining the first ray $[ps\,pa]$.

\paragraph{\tkzAttr{angle}{pb}}\label{ssub:attr_angle_pb}
A point defining the second ray $[ps\,pb]$.

\paragraph{\tkzAttr{angle}{raw}}\label{ssub:attr_angle_raw}
The oriented angle in radians as returned by \texttt{get\_angle\_(ps,pa,pb)}; it may be negative.

\paragraph{\tkzAttr{angle}{norm}}\label{ssub:attr_angle_norm}
The oriented angle normalised to the interval $[0,2\pi)$.

All values are numerical scalars and remain fixed once the object is created.

\subsection{Methods}

\vspace{1em}
\bgroup
\small
\captionof{table}{angle methods.}\label{angle:methods}
\begin{tabular}{ll}
\toprule
\textbf{Methods} & \textbf{Reference} \\
\midrule
\textbf{Creation} & \\
\midrule
\tkzMeth{angle}{angle(ps, pa, pb)} & [\ref{ssub:method_angle_new}] \\
\midrule
\textbf{Accessors} & \\
\midrule
\tkzMeth{angle}{get()} & [\ref{ssub:method_angle_get}] \\
\tkzMeth{angle}{value()} & [\ref{ssub:method_angle_value}] \\
\tkzMeth{angle}{deg()} & [\ref{ssub:method_angle_deg}] \\
\midrule
\textbf{Tests} & \\
\midrule
\tkzMeth{angle}{is\_direct()} & [\ref{ssub:method_angle_is_direct}] \\
\midrule
\textbf{Aliases} & \\
\midrule
\tkzMeth{angle}{to\_degrees()} & [\ref{ssub:method_angle_to_degrees}] \\
\bottomrule
\end{tabular}
\egroup


\subsubsection{\tkzMeth{angle}{get()}}
\label{ssub:method_angle_get}
Returns the three defining points:
\begin{verbatim}
local ps, pa, pb = alpha:get()
\end{verbatim}

\subsubsection{\tkzMeth{angle}{is\_direct()}}
\label{ssub:method_angle_is_direct}
Returns \verb|true| when the angle is positive (counterclockwise orientation).

\subsubsection{\tkzMeth{angle}{value()}}
\label{ssub:method_angle_value}
Returns the interior (non-oriented) angle in the range $[0,\pi]$:

\begin{tkzexample}[latex=.45\textwidth]
  \directlua{%
  init_elements()
  z.O = point(0, 1)
  z.T = point(2, 2)
  C.OT = circle(z.O, z.T)
  z.M = C.OT:point(.13)
  A.OTM = angle(z.O, z.T, z.M)
  tkzA = A.OTM:value()}
\begin{center}
  \begin{tikzpicture}
  \tkzGetNodes
  \tkzDrawLines(O,T O,M)
  \tkzDrawCircle(O,T)
  \tkzDrawPoints(O,T,M)
  \tkzLabelPoints(O,T,M)
  \tkzMarkAngle(T,O,M)
  \tkzLabelAngle[pos=1.5](T,O,M){%
   \tkzPN[3]{\tkzUseLua{tkzA}}}
\end{tikzpicture}
  \end{center}
\end{tkzexample}


\subsubsection{\tkzMeth{angle}{deg()}}
\label{ssub:method_angle_deg}
Returns the interior angle in degrees.

\subsubsection{\tkzMeth{angle}{to\_degrees()}}
\label{ssub:method_angle_to_degrees}
Alias of \texttt{deg()}.

\subsection{Example}

\begin{tkzexample}[latex=.35\textwidth]
\directlua{
  init_elements()
  z.A = point(0,0)
  z.B = point(3,0)
  z.C = point(1,2)
  A.alpha = angle(z.B, z.A, z.C)
  tex.print("Angle at A = "..A.alpha:deg().." degrees")
}
\end{tkzexample}

\medskip
This class is \emph{experimental} and may evolve in future versions.

\input{TKZdoc-elements-intersection.tex}
\input{TKZdoc-elements-constants.tex}
\input{TKZdoc-elements-tkz.tex}
\newpage
\section{Module utils}
The \tkzname{utils} module provides low-level numerical and parsing
utilities used internally by \tkzname{tkz-elements}.
These functions ensure safe numeric conversions and consistent
string formatting for TikZ output.


\subsection{Table of module functions \tkzname{utils}}

\begin{table}[htbp]
\centering
\caption{Functions of the module \tkzMod{utils}.}

\begin{tabular}{@{}ll@{}}
\toprule
\texttt{Function} & \texttt{Reference} \\
\midrule

\tkzFct{utils}{utils.parse\_point(str)}        & [\ref{sub:function_utils_parse_point}] \\

\tkzFct{utils}{utils.format\_number(r, n)}     & [\ref{sub:function_utils_format_number}]\\

\tkzFct{utils}{utils.format\_coord(x, decimals)} & [\ref{sub:function_utils_format_coord}]\\

\tkzFct{utils}{utils.format\_point(z, decimals)} & [\ref{sub:function_utils_format_point}]\\

\tkzFct{utils}{utils.checknumber(x, decimals)}  & [\ref{sub:function_utils_checknumber}]\\

\tkzFct{utils}{utils.almost\_equal(a, b, eps)}  & [\ref{sub:function_utils_almost_equal}]\\

\tkzFct{utils}{utils.wlog(...)}              & [\ref{sub:function_utils_wlog}]\\
\bottomrule
\end{tabular}
\end{table}



\subsection{Function \tkzFct{utils}{utils.parse\_point(str)}}
\label{sub:function_utils_parse_point}

Parses a string of the form \code{"(x,y)"} and returns the corresponding numeric coordinates. This function supports optional spaces and scientific notation.

\medskip
\texttt{Arguments: } \code{str} – A string representing a point, e.g., \code{"(3.5, -2.0)"}.

The function accepts:
\begin{itemize}
\item Decimal notation: \texttt{(3.5,-2)}
\item Scientific notation: \texttt{(1e-5,2E3)}
\item Optional spaces
\item Optional parentheses
\end{itemize}

\medskip
\texttt{Syntax: }

\begin{center}
\code{local x, y = utils.parse\_point("(1.5, -2.3)")}
\end{center}

\medskip
\texttt{Purpose: }

The function takes a string argument and parses it to extract the \code{x} and \code{y} components as numbers. The input string must follow the format \code{"(x, y)"} where \code{x} and \code{y} can be floating-point values written in decimal or scientific notation.


\medskip
\texttt{Returns: }

Two numbers: \texttt{x}, \texttt{y} – numeric coordinates as Lua numbers.
Two numerical values: the real and imaginary parts of the point.


\medskip
\texttt{Example usage: }

\begin{verbatim}
utils.parse_point("(3.5, -2)")
--  3.5   -2.0

utils.parse_point("(1e-5,2E3)")
-- 1e-05  2000

utils.parse_point(" -3.2e+1 , 4.5e-2 ")
-- -32  0.045
\end{verbatim}

Scientific and decimal forms are numerically equivalent:

\begin{verbatim}
local x1 = utils.parse_point("(1e-5,0)")
local x2 = utils.parse_point("(0.00001,0)")
-- x1 == x2  --> true
\end{verbatim}

\medskip
\texttt{Error handling}

Invalid input raises a \texttt{tex.error}:

\begin{verbatim}
utils.parse_point("(a,b)")
-- Invalid point string
\end{verbatim}



\subsection{Function \tkzFct{utils}{utils.format\_number(x, decimals)}}
\label{sub:function_utils_format_number}

This function formats a numeric value (or a numeric string) into a string representation with a fixed number of decimal places.

\texttt{Syntax: }

\begin{center}
\code{local str = utils.format\_number(math.pi, 3)}
\end{center}

\texttt{Purpose: }

The function converts a number (or a string that can be converted to a number) into a string with the specified number of decimal digits. It is especially useful when generating clean numerical output for display or export to \TIKZ{} coordinates.

\texttt{Arguments: }

\begin{itemize}
\item \code{x} – numeric value
\item \code{decimals} – Optional. number of decimal places (default: 5)
\end{itemize}

\texttt{Returns: } A string representing the value of \code{x} with the specified number of decimals.

\texttt{Features: }

\begin{itemize}
\item Automatically converts strings to numbers if possible.
\item Ensures consistent formatting for \TIKZ{} coordinates or LaTeX output.
\item Raises an error if the input is not valid.
\end{itemize}

\texttt{Example usage: }

\begin{verbatim}
utils.format_number(1e-5,5)
-- "0.00001"

utils.format_number(math.pi, 3)
-- "3.142"

utils.format_number("2.718281828", 2)
-- "2.72"

utils.format_number(2e3,2)
-- "2000.00"

utils.format_number(-4.2e-2,4)
-- "-0.0420"
\end{verbatim}


\subsection{Typical Use in tkz-elements}

Example of a complete parsing and formatting pipeline:

\begin{tkzexample}[latex=.4\textwidth]
\directlua{
init_elements()
local x,y = utils.parse_point("(1e-5,2e3)")

local sx = utils.format_number(x,5)
local sy = utils.format_number(y,2)
tex.print(sx)
tex.print('\\\\')
tex.print(sy)
}
\end{tkzexample}



\subsection{Function \tkzFct{utils}{utils.format\_coord(x, decimals)}}
\label{sub:function_utils_format_coord}

This function formats a numerical value into a string with a fixed number of decimal places. It is a lighter version of \tkzFct{utils}{format\_number}, intended for internal use when inputs are guaranteed to be numeric.

\texttt{Syntax: }

\begin{center}
\code{local s = utils.format\_coord(3.14159, 2)} \hfill $\rightarrow$ \code{"3.14"}
\end{center}

\texttt{Arguments: }

\begin{itemize}
\item \code{x} – A number (not validated).
\item \code{decimals} – Optional number of decimal places (default: 5).
\end{itemize}

\medskip
\texttt{Returns: }

A string with fixed decimal formatting.

\medskip
\texttt{Notes: }

This function is used internally by \tkzFct{path}{add\_pair\_to\_path} and other path-building methods.

Unlike \tkzFct{utils}{format\_number}, it does not perform input validation and should only be used with known numeric inputs.

\medskip
\texttt{Related functions: }

\begin{itemize}
\item \tkzFct{utils}{format\_number(x, decimals)} – safer alternative with validation
\end{itemize}

\medskip
\texttt{Precision considerations: }

Very small values may be rounded to zero depending on the number of decimals:

\begin{verbatim}
utils.format_number(1e-8,5)
-- "0.00000"

utils.format_number(1e-8,10)
-- "0.0000000100"
\end{verbatim}

\texttt{Important note: }

Lua may internally represent numbers using scientific notation:

\begin{verbatim}
print(1e-5)
-- 1e-05
\end{verbatim}

However, \verb|utils.format_number| always produces a fixed decimal
representation suitable for TikZ output.

\subsection{Interaction with \tkzname{path}}

The formatting utilities are especially important when building
TikZ paths from computed points.

\texttt{Example: Scientific notation inside a path}

Consider a very small coordinate obtained after a computation:

\begin{verbatim}
local x = 1e-5
local y = 2
\end{verbatim}

If inserted directly into a path without formatting:

\begin{verbatim}
PA.curve:add_point(point(x,y))
\end{verbatim}

Lua may internally represent the coordinate as:

\begin{verbatim}
(1e-05,2)
\end{verbatim}

While numerically correct, scientific notation is not ideal
for TikZ output.


\texttt{Complete example}

\begin{verbatim}
local p = path()
p:add_point(point(1e-5,0),5)
p:add_point(point(2e-5,1),5)

-- produces:
-- (0.00001,0.00000) -- (0.00002,1.00000)
\end{verbatim}





\subsection{Function \tkzFct{utils}{utils.checknumber(x, decimals)}}
\label{sub:function_utils_checknumber}

Validates and converts a number or numeric string into a fixed-format decimal string.

Alias of \verb|utils.format_number| (kept for backward compatibility)

\begin{verbatim}
utils.checknumber(1e-5,5)
-- "0.00001"
\end{verbatim}




\subsection{Function \tkzFct{utils}{utils.format\_point(z, decimals)}}
\label{sub:function_utils_format_point}

Converts a complex point into a string representation suitable for coordinate output.

\texttt{Syntax: }

\begin{center}
\code{local s = utils.format\_point(z, 4)} \hfill $\rightarrow$ \code{"(1.0000,2.0000)"}
\end{center}

\medskip
\texttt{Arguments: }

\begin{itemize}
\item \code{z} – A table with fields \code{re} and \code{im}.
\item \code{decimals} – Optional precision (default: 5).
\end{itemize}

\medskip
\texttt{Returns: }

A string representing the point as \code{"(x,y)"}.

\medskip
\texttt{Error handling.}

Raises an error if \code{z} does not have numeric \code{re} and \code{im} components.

\medskip
\texttt{Related functions: }

\begin{itemize}
\item \tkzFct{utils}{format\_coord}
\end{itemize}

\subsection{Function \tkzFct{utils}{utils.almost\_equal(a, b, epsilon)}}
\label{sub:function_utils_almost_equal}

Returns \code{true} if two numbers are approximately equal within a given tolerance.

\medskip
\texttt{Syntax: }

\begin{center}
\code{if utils.almost\_equal(x, y) then ... end}
\end{center}

\medskip
\texttt{Arguments: }

\begin{itemize}
\item \code{a}, \code{b} – Two numbers to compare.
\item \code{epsilon} – Optional tolerance (default: \code{tkz\_epsilon}).
\end{itemize}

\medskip
\texttt{Returns: }

A boolean: \code{true} if the values differ by less than the tolerance.



\subsection{Function \tkzFct{utils}{utils.wlog(...)}}
\label{sub:function_utils_wlog}

Logs a formatted message to the .log file only, with a \code{[tkz-elements]} prefix.

\medskip
\texttt{Syntax: }

\begin{center}
\code{utils.wlog("Internal value: \%s", tostring(value))}
\end{center}

\medskip
\texttt{Returns: }

No return value. Logging only.
\endinput

\input{TKZdoc-elements-maths.tex}

\part{Lua and Integration}


\input{TKZdoc-elements-lua.tex}
%--------------------------------------------------------------------
\newpage
\section{Transfers}
\label{sec:transfers}

The \tkzNamePack{tkz-elements} package strictly separates
\emph{mathematical computation} from \emph{graphical rendering}.
All geometric objects are created and manipulated internally in Lua,
while drawing is performed by \tkzNamePack{TikZ} or
\tkzNamePack{tkz-euclide}.

A \emph{transfer} is the explicit operation that communicates data
between the Lua computational layer and the \TeX\ rendering layer.

\medskip
Conceptually, the architecture may be summarized as:

\begin{center}
Computation (Lua) $\longrightarrow$ Transfer $\longrightarrow$ Rendering (TikZ/\TeX)
\end{center}

Objects remain invisible to \TeX\ until they are explicitly transferred.
This design guarantees numerical robustness and a clear separation
between computation and presentation.

%--------------------------------------------------------------------
\subsection{Conceptual Overview}

Three types of transfer mechanisms are available in
\tkzNamePack{tkz-elements}:

\begin{enumerate}
\item Immediate transfer via the \TeX\ input stream;
\item Object-based transfer through dedicated macros;
\item File-based transfer via external data files.
\end{enumerate}

Each mechanism serves a different purpose depending on the size
and nature of the data being communicated.

\subsection{Object Serialization via \texttt{tostring}}
\label{subsec:transfers_tostring}

The macro \tkzNameMacro{tkzUseLua} is the basic bridge from Lua to \TeX.
It evaluates a Lua expression and prints its result into the \TeX\ input stream.
Internally, it uses Lua's \texttt{tostring()} mechanism:

\begin{verbatim}
\def\tkzUseLua#1{\directlua{tex.sprint(#1)}}
\end{verbatim}

As a consequence, any \tkzNamePack{tkz-elements} object may define a Lua
\texttt{\_\_tostring} metamethod to control its textual \TeX\ representation.
This is used in particular for complex numbers (\tkzClass{point}).

\medskip
\texttt{Example with a point:}
\begin{tkzexample}[latex=.5\textwidth]
\directlua{
  init_elements()
  z.A = point(1, -2)
}
\tkzUseLua{tostring(z.A)}
\end{tkzexample}

Thus, the \emph{complex affix} associated with point $A$ is printed using
Lua’s \texttt{\_\_tostring} metamethod (via \texttt{tostring}).

\medskip
\texttt{Example with a matrix: }

The macro \tkzNameMacro{tkzUseLua} prints the \emph{value} of a Lua expression
as text (internally via \texttt{tostring(<expr>)}).%
\footnote{Hence objects may customize their inline representation by defining a Lua \texttt{\_\_tostring} metamethod.}


\medskip
\begin{tkzexample}[latex=.5\textwidth]
\directlua{
  init_elements()
  local a, b, c, d = 1, 2, 3, 4
  M.new = matrix({ { a, b }, { c, d } })}
  \tkzUseLua{M.new:print()}
\end{tkzexample}

\texttt{Other cases}
To transfer \code{nil}, \code{true}, and \code{false}, you must use \code{tostring}, However, you can directly transfer a number or a string of characters.

\begin{verbatim}
  \tkzUseLua{tostring(nil)}
  \tkzUseLua{tostring(true)}
  \tkzUseLua{"TANGENT"}
  \tkzUseLua{math.pi}
\end{verbatim}
%--------------------------------------------------------------------
\subsection{Immediate Transfer}
\label{subsec:transfer_immediate}

Immediate transfer occurs when Lua prints data directly into the
\TeX\ input stream using \texttt{tex.print}.

\subsubsection*{Numerical example}

\begin{tkzexample}[latex=.5\textwidth]
\directlua{
  init_elements()
  z.A = point(3,4)
  tex.print(point.abs(z.A))
}
\end{tkzexample}

The computed value is inserted directly into the document.

\subsubsection*{Boolean example}

\begin{tkzexample}[latex=.5\textwidth]

\directlua{
  init_elements()
  z.A = point(0, 0)
  z.B = point(3, 0)
  z.C = point(0, 2)
  T.ABC = triangle(z.A, z.B, z.C)
  tex.print(tostring(T.ABC:check_equilateral()))
}
\end{tkzexample}

This mechanism allows Lua computations to influence
\TeX\ logic through conditional statements.

\medskip
\texttt{Note: } It is possible to choose when to perform the transfer. A macro has been created for this purpose, but you can also create your own: it is called \tkzcname{tkzUseLua}. It is defined as follows:

\begin{mybox}
  \begin{verbatim}
  \def\tkzUseLua#1{\directlua{tex.print(#1)}}
\end{verbatim}
\end{mybox}

This macro prints the value of a Lua variable or expression directly into the TeX stream.

\medskip
\texttt{Example.}

 The following Lua code computes whether two lines intersect:

\directlua{
 z.a = point(4, 2)
 z.b = point(1, 1)
 z.c = point(2, 2)
 z.d = point(5, 1)
 L.ab = line(z.a, z.b)
 L.cd = line(z.c, z.d)
 det = (z.b - z.a) ^ (z.d - z.c)
 if det == 0 then bool = true
   else bool = false
 end
 x = intersection (L.ab, L.cd)}

The intersection of the two lines lies at
a point whose affix is: \tkzUseLua{x}

\vspace{1em}
\begin{minipage}{0.5\textwidth}
\begin{verbatim}
\directlua{
  z.a = point(4, 2)
  z.b = point(1, 1)
  z.c = point(2, 2)
  z.d = point(5, 1)
  L.ab = line(z.a, z.b)
  L.cd = line(z.c, z.d)
  det = (z.b - z.a) ^ (z.d - z.c)
  if det == 0 then bool = true
    else bool = false
  end
  x = intersection (L.ab, L.cd)}

The intersection of the two lines lies at
    a point whose affix is:\tkzUseLua{tostring(x)}


\begin{tikzpicture}
  \tkzGetNodes
  \tkzInit[xmin =0,ymin=0,xmax=5,ymax=3]
  \tkzGrid\tkzAxeX\tkzAxeY
  \tkzDrawPoints(a,...,d)
  \ifthenelse{\equal{\tkzUseLua{bool}}{%
  true}}{\tkzDrawSegments[red](a,b c,d)}{%
  \tkzDrawSegments[blue](a,b c,d)}
  \tkzLabelPoints(a,...,d)
\end{tikzpicture}
\end{verbatim}
\end{minipage}
\begin{minipage}{0.5\textwidth}
\begin{center}
\begin{tikzpicture}
  \tkzGetNodes
  \tkzInit[xmin =0,ymin=0,xmax=5,ymax=3]
  \tkzGrid\tkzAxeX\tkzAxeY
  \tkzDrawPoints(a,...,d)
  \ifthenelse{\equal{\tkzUseLua{bool}}{true}}{
  \tkzDrawSegments[red](a,b c,d)}{%
  \tkzDrawSegments[blue](a,b c,d)}
  \tkzLabelPoints(a,...,d)
\end{tikzpicture}
\end{center}
  \end{minipage}

%----------------------------------------------
% ---> INSERT additional custom examples here
%----------------------------------------------

%--------------------------------------------------------------------
\subsection{Object-Based Transfer}
\label{subsec:transfer_objects}

Geometric objects stored in Lua tables must be converted
into TikZ-compatible structures before drawing.

\subsubsection{Transferring points}
\label{transferring_points}

Points are internally stored as complex affixes.
The macro \tkzNameMacro{tkzGetNodes} converts them
into TikZ coordinates.

\begin{tkzexample}[latex=7cm]
  \directlua{
     init_elements()
     z.o   = point(0, 0)
     z.a_1 = point(2, 1)
     z.a_2 = point(1, 2)
     z.ap  = z.a_1 + z.a_2
     z.app = z.a_1 - z.a_2
  }

  \begin{center}
    \begin{tikzpicture}[ scale = 1.5]
       \tkzGetNodes
       \tkzDrawSegments(o,a_1 o,a_2 o,a' o,a'')
       \tkzDrawSegments[red](a_1,a' a_2,a')
       \tkzDrawSegments[blue](a_1,a'' a_2,a'')
       \tkzDrawPoints(a_1,a_2,a',o,a'')
       \tkzLabelPoints(o,a_1,a_2,a',a'')
    \end{tikzpicture}
  \end{center}
\end{tkzexample}

\subsubsection{Transferring paths and curves}


\begin{tkzexample}[latex=.4\textwidth]
\directlua{
 init_elements()
 PF.ex = pfct("(cos(t))^3", "(sin(t))^3")
 PA.curve = PF.ex:path(0, 2*math.pi, 100)
 }

\begin{center}
 \begin{tikzpicture}[scale=2]
  \tkzDrawCoordinates[smooth,red,fill=orange](PA.curve)
\end{tikzpicture}
 \end{center}
\end{tkzexample}

The Lua object is serialized into a coordinate list
interpreted by TikZ.

%----------------------------------------------
% ---> INSERT advanced fct / pfct examples here
%----------------------------------------------

%--------------------------------------------------------------------
\subsubsection{File-Based Transfer}


For large datasets or high-resolution curves,
it may be preferable to generate an external file.

This time, the transfer will be carried out using an external file. The following example is based on this one, but using a table.

\texttt{Generating a data file: } The file \texttt{tmp.table} is created using the $f$ function.

The file \texttt{tmp.table} now contains numerical data.
 \texttt{Using the generated file: } Here, the curve is plotted using \tkzNamePack{TikZ}.



\begin{tkzexample}[latex=.5\textwidth]
 \directlua{
  init_elements()
   z.a = point(1, 0)
   z.b = point(3, 2)
   z.c = point(0, 2)
   A,B,C = tkz.parabola (z.a, z.b, z.c)

 function f(t0, t1, n)
  local out=assert(io.open("tmp.table","w"))
  local y
  for t = t0,t1,(t1-t0)/n  do
   y = A*t^2+B*t +C
   out:write(utils.checknumber(t), " ",
     utils.checknumber(y), " i\string\n")
  end
  out:close()
 end}

\begin{tikzpicture}
   \tkzGetNodes
   \tkzInit[xmin=-1,xmax=5,ymin=0,ymax=5]
   \tkzDrawX\tkzDrawY
   \tkzDrawPoints[red,size=2](a,b,c)
   \directlua{f(-1,3,100)}%
   \draw[domain=-1:3] plot[smooth]
        file {tmp.table};
\end{tikzpicture}
 \end{tkzexample}





File-based transfer provides scalability and allows
interoperability with external plotting tools such as
\texttt{pgfplots}.

%----------------------------------------------
% ---> INSERT high-resolution example here
%----------------------------------------------

%--------------------------------------------------------------------
\subsection{Dynamic \TeX--Lua Interaction}
\label{subsec:dynamic_transfer}

Transfers are bidirectional.
\TeX\ may pass parameters to Lua for computation.

\medskip
\begin{verbatim}
\newcommand{\ComputeSquare}[1]{%
  \directlua{tex.print(#1 * #1)}%
}
\end{verbatim}

\medskip
Usage:
\newcommand{\ComputeSquare}[1]{%
  \directlua{tex.print(#1 * #1)}%
}
\begin{verbatim}
The square of 5 is \ComputeSquare{5}.
\end{verbatim}

\medskip
The square of 5 is \ComputeSquare{5}.

\medskip

This interaction enables dynamic parameterized geometry,
where document input directly controls Lua computations.

%----------------------------------------------
% ---> INSERT dynamic geometry example here
%----------------------------------------------

\medskip
Transfers in \tkzNamePack{tkz-elements} may therefore be
immediate, object-based, or file-based,
depending on the size and structure of the data.
This layered architecture ensures a clean separation
between computation and graphical representation.
%====================================================================
\section{TeX Interface Macros (\texttt{tkz-elements.sty})}
\label{sec:sty_interface}

\subsection{Purpose}

The file \texttt{tkz-elements.sty} provides a small set of macros that
act as a bridge between \TeX{} and the Lua engine used by
\tkzNamePack{tkz-elements}. Their main goals are:

\begin{itemize}
\item printing Lua values safely in \TeX{};
\item transferring geometric objects (points, paths, circles) from Lua to TikZ;
\item simplifying the drawing stage when computations are done in Lua.
\end{itemize}

\medskip
Unless stated otherwise, these macros assume that \tkzNamePack{tkz-elements}
has been initialized in Lua (typically via |\directlua{init_elements()}|.

%--------------------------------------------------------------------
\subsection{Summary of provided macros}


\begin{tabular}{ll}
\toprule
\texttt{Macro} & \texttt{Reference} \\
\midrule
\tkzMacro{tkz-elements}{tkzUseLua\{...\}}               & \\
\tkzMacro{tkz-elements}{tkzPrintNumber\{...\}}          & \\
\tkzMacro{tkz-elements}{tkzEraseLuaObj\{name\}}         & \\
\tkzMacro{tkz-elements}{tkzDrawCoordinates(\dots)}      & \\
\tkzMacro{tkz-elements}{tkzGetPointsFromPath\{P\}\{A\}} & \\
\tkzMacro{tkz-elements}{tkzDrawPointsFromPath(\dots)}   & \\
\tkzMacro{tkz-elements}{tkzDrawSegmentsFromPaths(\dots)}& \\
\tkzMacro{tkz-elements}{tkzDrawCirclesFromPaths(\dots)} & \\
\bottomrule
\end{tabular}

\medskip
\texttt{Note:} The exact internal Lua object layout may evolve; these macros
are designed to remain stable at the user level.

%====================================================================
\subsection{Macro \tkzMacro{tkz-elements}{tkzUseLua}}
\label{subsec:tkzUseLua}

\medskip
\texttt{Syntax: }

\verb|tkzUseLua{<lua-expression>\}|


\medskip
\texttt{Description: }
Evaluates the Lua expression and prints its result in the \TeX{} stream.

\medskip
\texttt{Argument: }
\begin{itemize}
\item \tkzname{<lua-expression>}: a Lua expression returning a value.
\end{itemize}

\medskip
\texttt{Notes: }
\begin{itemize}
\item This macro is meant for \emph{values} (numbers, booleans, strings).
\item For Lua objects (point, line, circle, path, \dots), prefer dedicated
transfer/draw macros or call an explicit Lua method that returns a value.
\end{itemize}

\texttt{Example: }
\begin{minipage}{.52\textwidth}
\begin{verbatim}
\directlua{init_elements()
  z.A = point(0,0)
  z.B = point(4,2)
}
The slope is \tkzUseLua{(z.B.im-z.A.im)/(z.B.re-z.A.re)}.
\end{verbatim}
\end{minipage}

%%%%%%%%%%%%%%%%%%%%%%%%%%%%%%%

This macro evaluates a Lua expression and inserts its result in the document.


\texttt{Syntax: }\verb|\tkzUseLua{<lua-code>}|

\texttt{Example usage: }

\begin{mybox}
\begin{verbatim}
\tkzUseLua{checknumber(math.pi)}
\end{verbatim}
\end{mybox}


%====================================================================
\subsection{Macro \tkzMacro{tkz-elements}{tkzPrintNumber}}
This macro formats and prints a number using PGF’s fixed-point output, with a default precision of 2 decimal places.

\medskip

%\texttt{Syntax: }\verb|\tkzPrintNumber[<precision>]{<number>}|
\texttt{Syntax: }
\begin{quote}
\tkzMacro{tkz-elements}{tkzPrintNumber\{<lua-expression>\}} \qquad
\tkzMacro{tkz-elements}{tkzPN\{<lua-expression>\}}
\end{quote}

\texttt{Example: }

\begin{mybox}
\begin{verbatim}
\tkzPrintNumber{pi} % outputs 3.14
\tkzPrintNumber[4]{sqrt(2)} % outputs 1.4142
\end{verbatim}
\end{mybox}

\texttt{Alias: } The macro \tkzMacro{tkz-elements}{tkzPN} is a shorthand for \tkzMacro{tkz-elements}{tkzPrintNumber}.


\texttt{Example: }
\begin{minipage}{.52\textwidth}
\begin{verbatim}
\directlua{init_elements()
  x = math.pi/7
}
$\alpha = \tkzPN{x}$.
\end{verbatim}
\end{minipage}


%====================================================================
\subsection{Macro \tkzMacro{tkz-elements}{tkzEraseLuaObj}}

\label{sub:macro_tkzEraseLuaObj}

\texttt{Syntax: }
\begin{quote}
\tkzMacro{tkz-elements}{tkzEraseLuaObj\{<name>\}}
\end{quote}



\texttt{Description: }
Removes a stored Lua object from the global storage table (typically \code{z},
\code{L}, \code{C}, \code{CO}, etc., depending on your conventions).


\texttt{Syntax: } \verb|\tkzEraseLuaObj{<object>}|

\texttt{Example usage: }

\begin{mybox}
\begin{verbatim}
\tkzEraseLuaObj{z.A}
\tkzEraseLuaObj{T.ABC}
\end{verbatim}
\end{mybox}


\texttt{Notes: }
Use this macro to avoid name collisions when compiling large documents
or when reusing the same object names across figures.



%====================================================================
\subsection{Macro \tkzMacro{tkz-elements}{tkzPathCount}}
This macro retrieves the number of vertices in a Lua path and stores the result in a TeX macro.

\texttt{Syntax: }\verb|\tkzPathCount(<lua-path>){<MacroName>}|

\texttt{Arguments: }

\verb|<lua-path>|: identifier of a Lua \verb|path| object.
\verb|<MacroName>|: name of a TeX macro (without the backslash) that will store the count.


\texttt{Example usage: }

\begin{tkzexample}[latex=.5\textwidth]
\directlua{
 z.A = point(0, 0)
 z.B = point(5, 4)
 L.AB = line(z.A, z.B)
 PA.AB = L.AB:path(3)
 z.O = point(2, 0)
 }
\tkzPathCount(PA.AB){N}
\begin{center}
\begin{tikzpicture}
\tkzGetNodes
\tkzDrawCoordinates[blue](PA.AB)
\tkzDrawPointsFromPath[red,size=2](PA.AB)
  \foreach \i in {1,...,\N}{
 \tkzGetPointFromPath(PA.AB,\i){P\i}
 \tkzDrawSegment[orange](O,P\i)
  }
\end{tikzpicture}
\end{center}
\end{tkzexample}



%====================================================================
\subsection{Macro \tkzMacro{tkz-elements}{tkzDrawCoordinates}}
This macro draws a curve using the TikZ \code{plot coordinates} syntax, with coordinates generated by Lua.


\texttt{Syntax: } |\tkzDrawCoordinates[<options>](lua-path)|

\texttt{Example usage: }

%\begin{mybox}
\begin{tkzexample}[latex=.5\textwidth]
\directlua{
  z.A = point(0, 0)
  z.B = point(6, 0)
  z.C = point(1, 4)
  T.ABC = triangle(z.A, z.B, z.C)
  z.M = T.ABC.bc.mid
  z.N = T.ABC.ca.mid
  PA.path = T.ABC:path(z.M,z.N,4)
}
\begin{tikzpicture}[gridded]
\tkzGetNodes
\tkzDrawPolygon(A,B,C)
\tkzDrawCoordinates[orange,thick](PA.path)
\tkzDrawPoints(A,B,C,M,N)
\tkzLabelPoints(A,B)
\tkzLabelPoints[above](C)
\end{tikzpicture}
\end{tkzexample}
%\end{mybox}

\medskip
\texttt{Description: }
Draws a Lua \code{path} object as a TikZ path. This macro is typically used
after the computational phase in Lua. Draws the polyline defined by the sequence of points in a Lua path, using explicit \code{(x,y)} coordinates.


\medskip
\texttt{Alias: } % \let\tkzDrawPath\tkzDrawCoordinates
\label{subsec:tkzDrawPath}

\begin{quote}
\tkzMacro{tkz-elements}{tkzDrawPath[<tikz-options>](lua-path)}
\end{quote}

%====================================================================

\subsection{Macro \tkzMacro{tkz-elements}{tkzDrawPointsFromPath}}
\label{sub:macro_tkzDrawPointsFromPath}

This macro draws all points of a path using \tkzname{tkzDrawPoint}, without exporting point names.

\medskip
\texttt{Syntax: }
\verb|\tkzDrawPointsFromPath[<options>](lua-path)|

\medskip
\texttt{Example usage: }

\begin{tkzexample}[latex=.5\textwidth]
\directlua{
init_elements()
LP.test = list_point()
LP.test:add(point(0, 0))
LP.test:add(point(1, 0))
LP.test:add(point(1, 1))
PA.curve = LP.test:as_path()
}
\begin{tikzpicture}
\tkzDrawCoordinates[blue](PA.curve)
\tkzDrawPointsFromPath[red,size=2](PA.curve)
\end{tikzpicture}
\end{tkzexample}

%====================================================================


%====================================================================
\subsection{Macro \tkzMacro{tkz-elements}{tkzGetPointsFromPath}}
This macro defines TikZ points from a Lua path. Each point is named with a base name followed by an index.

\texttt{Syntax: }\verb|\tkzGetPointsFromPath[<options>](<path>,<basename>)|

\medskip
\texttt{Description: }
Exports the points of a Lua path as named \TeX/TikZ points.
The naming scheme is \code{<prefix>1}, \code{<prefix>2}, \dots

\medskip
\texttt{Example: }

\begin{tkzexample}[latex=.5\textwidth]
 \directlua{
  init_elements()
  PA.g = path({ "(0, 0)", "(1, 0)","(1, 1)" })
}
\begin{tikzpicture}
  \tkzGetPointsFromPath(PA.g,A)
  \tkzDrawPolygon(A_1,A_2,A_3)
  \tkzDrawPoints(A_1,A_2,A_3)
\end{tikzpicture}
\end{tkzexample}

%====================================================================




\subsection{Macro \tkzMacro{tkz-elements}{tkzGetPointFromPath}}
This macro extracts the i-th vertex of a Lua path and defines it as a TikZ point with a given name.

\texttt{Syntax: }\verb|\tkzGetPointFromPath(<PathLua>,<index>){<PointName>}|

\texttt{Arguments: }

\verb|<PathLua>|: identifier of a Lua \verb|path| object.
\verb|<index>|: index of the vertex to retrieve (TeX number or numeric macro).
\verb|<PointName>|: name of the TikZ point to create.
\medskip


\texttt{Example usage:}

\begin{mybox}
\begin{verbatim}
% Get the 3rd vertex of PA.A and name it P3
\tkzGetPointFromPath(PA.A,3){P3}
% With counting and a loop
\tkzPathCount(PA.A){N}
\foreach \i in {1,...,\N}{
\expandafter\tkzGetPointFromPath\expandafter(PA.A,\i){P\i}
}
\end{verbatim}
\end{mybox}
%====================================================================




\subsection{Macro \tkzMacro{tkz-elements}{tkzDrawSegmentsFromPaths}}
\label{sub:macro_tkzDrawSegmentsFromPaths}

\texttt{Description: }
Draws segments encoded in one or more Lua paths (pairs of points or consecutive
points, depending on the chosen convention).

\texttt{Syntax: }\verb|\tkzDrawSegmentsFromPaths[<options>](<lua-path-start>,<lua-path-end>)|

\texttt{Example: }

\begin{tkzexample}[vbox]
\directlua{
 init_elements()
 z.A = point(0, 0)
 z.B = point(3, 0)
 C.AB = circle(z.A, z.B)
 PA.c = path()
 PA.u = path()
 PA.v = path()
 for i = 0, 18, 2 do
   local angle = i * 2 * math.pi / 18
   local p = point(polar(8, angle))
   local Lu, Lv = C.AB:tangent_from(p)
   local u , v  = Lu.pb, Lv.pb
   PA.c:add_point(p)
   PA.u:add_point(u)
   PA.v:add_point(v)
 end}

\begin{center}
 \begin{tikzpicture}[scale =.6]
  \tkzGetNodes
  \tkzDrawSegmentsFromPaths[draw,red](PA.c,PA.u)
  \tkzDrawSegmentsFromPaths[draw,red](PA.c,PA.v)
  \tkzDrawCircle[blue](A,B)
  \tkzDrawPoints(A,I,u,v)
\end{tikzpicture}
 \end{center}
\end{tkzexample}
%====================================================================
\subsection{Macro \tkzMacro{tkz-elements}{tkzDrawCirclesFromPaths}}
\label{sub:macro_tkzDrawCirclesFromPaths}
This macro draws circles defined by two paths: one for the centers and one for the points through which each circle passes.

\medskip
\texttt{Syntax: }\verb|\tkzDrawCirclesFromPaths[<options>](<PA.center>,<PA.through>)|

\medskip
\texttt{Example usage: }

\begin{mybox}
\begin{verbatim}
\tkzDrawCirclesFromPaths[blue](PA.O, PA.A)
\end{verbatim}
\end{mybox}

See the document \tkzname{Euclidean Geometry} presented in \href{http://altermundus.fr}{altermundus.fr}for other examples.

\medskip
\begin{tkzexample}[vbox]
\directlua{
 init_elements()
 z.A = point(-4, -4)
 z.B = point(4, -4)
 L.AB = line(z.A, z.B)
 C.AB = circle(z.A, z.B)
 S.AB = L.AB:square()
 _, _, z.C, z.D = S.AB:get()
 z.O = S.AB.ac.mid
 z.a = S.AB.ab.mid
 z.c = S.AB.cd.mid
 z.o1 = tkz.midpoint(z.a, z.O)
 z.o = tkz.midpoint(z.c, z.O)
 z.b = S.AB.bc.mid
 z.g = (z.o1 - z.O) + (z.b - z.O)
 z.o2 = intersection(line(z.c, z.g), line(z.O,z.b))
 z.o3 = (z.o - z.O) + (z.o2 - z.O)
 z.j = intersection(line(z.c, z.b), line(z.o2,z.o3))
 local C2 = circle(z.O, z.a)
 local C1 = circle(z.o, z.c)
 local C3 = circle(z.o3, z.j)
 PA.pc, PA.pt, n = C1:CCC(C2, C3)
}

\begin{center}
\begin{tikzpicture}[scale=1]
\tkzGetNodes
\tkzDrawPolygon[cyan](A,B,C,D)
\tkzDrawCircles(O,a o,c o1,a o2,b o3,j)
\tkzDrawCirclesFromPaths[draw,red](PA.pc,PA.pt)
\end{tikzpicture}
\end{center}
\end{tkzexample}


%===============================================================

\subsection{Macro \tkzMacro{tkz-elements}{tkzDrawFromPointToPath}}
This macro draws a set of line segments from a given TikZ point to every point of a Lua \emph{path} object.

\medskip
\texttt{Syntax: }\verb|\tkzDrawFromPointToPath[<options>](<point-tikz>,<lua-path>)|

\medskip
\texttt{Arguments: }

\verb|<options>| (optional): pgf/tikz drawing options (color, thickness, etc.).

\verb|<point-tikz>|: the name of a TikZ point (already defined).

\verb|<lua-path>|: identifier of a Lua \verb|path| object.

\medskip
\texttt{Example usage: }

\medskip
\begin{tkzexample}[vbox]
\directlua{
 init_elements()
 z.O = point(0, 0)
 z.A = point(5, 1)
 z.Ta = point(8, 2)
 C.A = circle(z.A, z.Ta)
 L.OA = line(z.O,z.A)
 z.x, z.y = intersection(C.A, L.OA,{near=z.O})
 C.A.through = z.y
 z.m = tkz.midpoint(z.O, z.x)
 z.n = tkz.midpoint(z.O, z.y)
 z.w = tkz.midpoint(z.m, z.n)
 PA.A = path()
 PA.c = path()
 for t = 0, 1 - 1e-12, 0.1 do
   PA.A:add_point(C.A:point(t))
   local m = tkz.midpoint(z.O, C.A:point(t))
   PA.c:add_point(m)
 end
}
 \begin{center}
  \begin{tikzpicture}[scale=.8]
   \tkzGetNodes
   \tkzDrawCircle(A,Ta)
   \tkzDrawLine(O,A)
   \tkzDrawPoints(O,A,m,n,w,x,y)
   \tkzLabelPoints(O,A,m,n,w)
   \tkzDrawCircle[red](w,n)
   \tkzDrawPointsFromPath[red](PA.c)
   \tkzDrawPointsFromPath[blue](PA.A)
   \tkzDrawSegmentsFromPaths[draw,teal](PA.c,PA.A)
   \tkzDrawFromPointToPath[orange](O,PA.A)
 \end{tikzpicture}
  \end{center}
\end{tkzexample}



\subsection{Archimedes spiral, a complete example}

This example is built with the help of some of the macros presented in this section. It is inspired by an example created with tkz-euclide by Jean-Marc Desbonnez.

\vspace{1em}
\begin{tkzexample}[vbox]
\directlua{
  z.O = point(0, 0)
  z.A = point(5, 0)
  L.OA = line(z.O, z.A)
  PA.spiral = path()
  PA.center = path()
  PA.through = path()
  PA.ray = path()
  for i = 1, 48 do
      local k = i / 48
      z["P"..i] = L.OA:point(k)
      local angle = i * tkz.tau / 48
      z["r"..i] = point(polar(5, angle))
      ray = line(z.O, z["r"..i])
      circ = circle(z.O, z["P"..i])
      z["X"..i], _ = intersection(ray, circ, { near = z["r"..i] })
      PA.spiral:add_point(z["X"..i])
      PA.center:add_point(z.O)
      PA.through:add_point(z["P"..i])
      PA.ray:add_point(z["r"..i])
  end}
\begin{center}
  \begin{tikzpicture}[scale = .8]
  \tkzGetNodes
  \tkzDrawCircle(O,A)
  \tkzDrawCirclesFromPaths[draw,
      orange!50](PA.center,PA.through)
  \tkzDrawSegmentsFromPaths[draw,teal](PA.center,PA.ray)
  \tkzDrawPath[red,thick,smooth](PA.spiral)
  \tkzDrawPointsFromPath[blue, size=2](PA.spiral)
  \end{tikzpicture}
\end{center}

\end{tkzexample}



\endinput


%--------------------------------------------------------------------
\newpage
\section{Class \tkzClass{fct}}
\label{sec:class_fct}

The \tkzClass{fct} class represents a real-valued function
\[
x \longmapsto f(x),
\]
defined and evaluated in Lua. A \tkzClass{fct} object encapsulates either
(i) a numerical expression (given as a string in the variable \texttt{x}),
or (ii) a callable Lua function, together with methods for evaluation,
sampling, and geometric construction.

By convention, function objects are stored in a Lua table named
\tkzVar{fct}{F}. Although this name is not mandatory, using \texttt{F} is
strongly recommended for consistency across examples and documentation.
Each entry of \texttt{F} associates a symbolic name (such as \texttt{fa})
with an object of class \tkzClass{fct}. If a custom table name is used, it
must be initialized manually. The \tkzFct{tkz-elements}{init\_elements()}
function will reset the \texttt{F} table if it has already been defined.

\paragraph{Mathematical expressions.}
All expressions defining a \tkzClass{fct} object are evaluated using standard
Lua rules. A \tkzClass{fct} object can be evaluated at a real number, sampled on an
interval, converted into a \tkzClass{path} (for drawing with TikZ), or exported
to a data file.

\paragraph{Important.}
In function and parametric function definitions
(\tkzClass{fct} and \tkzClass{pfct}),
expressions must be written using standard mathematical notation.
The Lua prefix \texttt{math.} must \emph{not} be used.

In contrast, in all other methods and Lua code fragments of
\tkzNamePack{tkz-elements},
standard Lua syntax applies, and the use of the prefix \texttt{math.}
is required whenever a Lua mathematical function is called. For convenience, users may define local aliases in their Lua code, such as
\code{local sin, cos, pi = math.sin, math.cos, math.pi}.


\medskip
Correct:
\begin{verbatim}
sin(x), exp(-x^2)
\end{verbatim}

Incorrect (in \tkzClass{fct} and \tkzClass{pfct}):
\begin{verbatim}
math.sin(x), math.exp(-x^2)
\end{verbatim}

\subsection{Methods of the class fct}
The \tkzClass{fct} class offers a compact set of methods, centered around:
\emph{(i)} creating a function object, \emph{(ii)} evaluating it, and
\emph{(iii)} producing a \tkzClass{path} or a data file.

\vspace{1em}

\bgroup
  \small
  \captionof{table}{fct methods.}\label{fct:methods}
  \begin{tabular}{ll}
  \toprule
  \textbf{Methods} & \texttt{Reference} \\
  \midrule
  \textbf{Creation} & \\
  \midrule
  \tkzFct{fct}{new(expr\_or\_fn)}              & [\ref{ssub:method_fct_new}] \\
  \tkzFct{fct}{compile(expr)}                 & [\ref{ssub:method_fct_compile}] \\
  \midrule
  \textbf{Reals / evaluation} & \\
  \midrule
  \tkzMeth{fct}{eval(x)}                     & [\ref{ssub:method_fct_value}] \\
  \midrule
  \textbf{Points / Paths} & \\
  \midrule
  \tkzMeth{fct}{point(x)}                     & [\ref{ssub:method_fct_point}] \\
  \tkzMeth{fct}{path(xmin,xmax,n)}            & [\ref{ssub:method_fct_path}] \\
  \bottomrule
  \end{tabular}
\egroup

%-------------------- details ---------------------------------------
\subsubsection{Constructor \tkzFct{fct}{new(expr\_or\_fn)}}
\label{ssub:method_fct_new}

\texttt{Description: }Creates a function object from either a Lua expression
(string) or a callable Lua function. When a string is provided, it represents
an expression in the variable \texttt{x}, evaluated using standard Lua rules
(therefore requiring \texttt{math.} prefixes).

\texttt{Arguments: }\par
\begin{itemize}
\item \code{expr\_or\_fn}: either
  \begin{itemize}
    \item a string representing an expression in \texttt{x}, or
    \item a callable Lua function \texttt{f(x)}.
  \end{itemize}
\end{itemize}

\texttt{Returns: } a \tkzClass{fct} object.

\texttt{Example:}
\begin{verbatim}
\directlua{
  init_elements()
  F.f = fct:new("x*math.exp(-x^2)+1")
  tex.print(F.f:value(0))
}
\end{verbatim}

\subsubsection{Function \tkzFct{fct}{compile(expr)}}
\label{ssub:method_fct_compile}

\texttt{Description: }Compiles an expression and returns a callable Lua function
(or wraps it into a \tkzClass{fct}, depending on the implementation). This is a
low-level helper mainly intended for internal use.

\texttt{Argument: } \code{expr} (string, expression in variable \texttt{x}). \qquad
\texttt{Returns: } a callable function or a \tkzClass{fct} object (depending on implementation).

\subsubsection{Method \tkzMeth{fct}{eval(x)}}
\label{ssub:method_fct_value}

\texttt{Description: }Evaluates the function at the real number \code{x}.

\texttt{Returns: } a number (which may be \texttt{nan} or infinite if the expression is not defined).

\begin{verbatim}
  init_elements()
  F.fa = fct("sin(x) + x")
  PA.curve = F.fa:path(-2, 5, 200)
  z.A = F.fa:point(1.3)
  tex.print(F.fa:eval(math.pi/2))
\end{verbatim}


\subsubsection{Method \tkzMeth{fct}{point(x)}}
\label{ssub:method_fct_point}

\texttt{Description: }Constructs the point $(x,f(x))$ associated with the function.
This method provides a direct bridge between numerical evaluation and geometric
construction.

\texttt{Returns: } a \tkzClass{point}.

\texttt{Example: }
\begin{verbatim}
  init_elements()
  F.fa = fct("sin(x)+x")
  PA.curve = F.fa:path(-2, 5, 200)
  z.A = F.fa:point(1.5)
  % or z.A = point(1.3, F.fa:eval(1.3) )
\end{verbatim}

\subsubsection{Method \tkzMeth{fct}{path(xmin,xmax,n)}}
\label{ssub:method_fct_path}

\texttt{Description: }Samples the function on $[xmin,xmax]$ using \code{n}
subdivisions and returns a \tkzClass{path}. Invalid values are skipped (depending
on the internal policy).

\texttt{Returns: } a \tkzClass{path}.


\texttt{Example: }

\begin{tkzexample}[latex=.5\textwidth]
\directlua{
  init_elements()
  F.fa = fct("sin(x) + x")
  PA.curve = F.fa:path(-1, 5, 200)
  z.A = F.fa:point(math.pi/2)
}
\begin{tikzpicture}[scale=.75,gridded]
  \tkzInit[xmin=-1,xmax=5,ymin=-1,ymax=2]
  \tkzDrawX\tkzDrawY
  \tkzGetNodes
  \tkzDrawCoordinates[smooth,blue,thick](PA.curve)
  \tkzDrawPoint[red](A)
  \tkzDrawPointsOnGraph[blue]{0,1,2}{fa}
  \tkzDrawPointOnGraph[red]{-1}{fa}
\end{tikzpicture}
\end{tkzexample}

%--------------------------------------------------------------------
\subsection{Macros \tkzMacro{tkz-elements}{tkzDrawPointOnGraph} and
\tkzMacro{tkz-elements}{tkzDrawPointsOnGraph}}

These macros allow drawing points on the graph of a function defined
in the module system.

The function must already exist in the module table \texttt{F}.

%--------------------------------------------------------------------

\subsubsection{Macro \tkzMacro{tkz-elements}{tkzDrawPointOnGraph}}

\paragraph{Syntax}
\begin{verbatim}
\tkzDrawPointOnGraph[<TikZ options>]{<x>}{<name>}
\end{verbatim}

\paragraph{Description}
This macro draws a point belonging to the graph of a real function stored
in the module \texttt{F}.

The drawn point has coordinates
\[
(x,\; f(x)),
\]
where \texttt{f} denotes the function object \texttt{F.<name>}.

\paragraph{Arguments}
\begin{itemize}
  \item \texttt{<TikZ options>} (optional): graphical options applied to the point
  \item \texttt{<x>}: abscissa of the point
  \item \texttt{<name>}: name of the function stored in \texttt{F}
\end{itemize}

\paragraph{Example}
\begin{verbatim}
\tkzDrawPointOnGraph[blue]{4}{fa}
\end{verbatim}

%--------------------------------------------------------------------
\subsubsection{Macro \tkzMacro{tkz-elements}{tkzDrawPointsOnGraph}}

\paragraph{Syntax}
\begin{verbatim}
\tkzDrawPointsOnGraph[<TikZ options>]{<x1,x2,...,xn>}{<name>}
\end{verbatim}

\paragraph{Description}
This macro draws several points belonging to the graph of a real function
stored in the module \texttt{fct}.

For each value \(x_i\) in the list, a point of coordinates
\[
(x_i,\; f(x_i))
\]
is computed and drawn.

\paragraph{Arguments}
\begin{itemize}
  \item \texttt{<TikZ options>} (optional): graphical options applied to all points
  \item \texttt{<x1,x2,...,xn>}: comma-separated list of abscissas
  \item \texttt{<name>}: name of the function stored in \texttt{F}
\end{itemize}

\paragraph{Example}
\begin{verbatim}
\tkzDrawPointsOnGraph[red]{-2,0,2,3,5}{fa}
\end{verbatim}

\paragraph{Remarks}
\begin{itemize}
  \item These macros assume that the function is already defined in the module
        \texttt{F}.
  \item All evaluations are performed in Lua.
  \item The macros are compatible with paths drawn using
        \texttt{\textbackslash tkzDrawCoordinates}.
\end{itemize}

\texttt{Example: }
\begin{tkzexample}[latex=.4\textwidth]
\directlua{
  init_elements()
  F.fa = fct("x*exp(-x^2)+1")
  PA.curve = F.fa:path(-3, 5, 20)
}
\begin{tikzpicture}[gridded]
  \tkzInit[xmin=-3,xmax=5,ymin=-1,ymax=2]
  \tkzDrawX\tkzDrawY
  \tkzDrawCoordinates[smooth,cyan](PA.curve)
  \tkzDrawPointsOnGraph[red]{-2,0,2,3,5}{fa}
  \tkzDrawPointOnGraph[blue]{1}{fa}
\end{tikzpicture}
\end{tkzexample}
\endinput
%--------------------------------------------------------------------
\newpage
\section{Class \tkzClass{pfct}}
\label{sec:class_pfct}

The \tkzClass{pfct} class represents a parametric curve
\[
t \longmapsto \bigl(x(t),y(t)\bigr),
\]
defined and evaluated in Lua. A \tkzClass{pfct} object encapsulates two
expressions (or callable Lua functions) describing the $x$- and $y$-components
of the curve, together with methods for evaluation, sampling, and geometric
construction.

By convention, parametric function objects are stored in a Lua table named
\tkzVar{pfct}{PF}. Although this name is not mandatory, using \texttt{PF} is
strongly recommended for consistency across examples and documentation. Each
entry of \texttt{PF} associates a symbolic name (such as \texttt{lis}) with an
object of class \tkzClass{pfct}. If a custom table name is used, it must be
initialized manually. The \tkzFct{tkz-elements}{init\_elements()} function will
reset the \texttt{PF} table if it has already been defined.

\paragraph{Mathematical expressions.}
All expressions defining a \tkzClass{pfct} object are evaluated using standard
Lua rules. A \tkzClass{pfct} object can be evaluated at a parameter value, converted into
points, sampled into a \tkzClass{path}, or exported to a data file suitable for
TikZ \texttt{plot file} input.

\paragraph{Important.}
In function and parametric function definitions
(\tkzClass{fct} and \tkzClass{pfct}),
expressions must be written using standard mathematical notation.
The Lua prefix \texttt{math.} must \emph{not} be used.

In contrast, in all other methods and Lua code fragments of
\tkzNamePack{tkz-elements},
standard Lua syntax applies, and the use of the prefix \texttt{math.}
is required whenever a Lua mathematical function is called. For convenience, users may define local aliases in their Lua code, such as
\code{local sin, cos, pi = math.sin, math.cos, math.pi}.

\medskip
Correct:
\begin{verbatim}
sin(x), exp(-x^2)
\end{verbatim}

Incorrect (in \tkzClass{fct} and \tkzClass{pfct}):
\begin{verbatim}
math.sin(x), math.exp(-x^2)
\end{verbatim}

\subsection{Methods of the class pfct}
\vspace{1em}

\bgroup
  \small
  \captionof{table}{pfct methods.}\label{pfct:methods}
  \begin{tabular}{ll}
  \toprule
  \texttt{Methods} & \texttt{Reference} \\
  \midrule
  \texttt{Constructor} & \\
  \midrule
  \tkzFct{pfct}{new(exprx,expry)}       & [\ref{ssub:method_pfct_new}] \\
  \tkzFct{pfct}{compile(exprx,expry)}   & [\ref{ssub:method_pfct_compile}] \\
  \midrule
  \texttt{Methods Returning a Real Number } & \\
  \midrule
  \tkzMeth{pfct}{x(t)}                 & [\ref{ssub:method_pfct_x}] \\
  \tkzMeth{pfct}{y(t)}                 & [\ref{ssub:method_pfct_y}] \\
  \tkzMeth{pfct}{point(t)}             & [\ref{ssub:method_pfct_point}] \\
  \midrule
  \texttt{Methods Returning a Path} & \\
  \midrule
  \tkzMeth{pfct}{path(tmin,tmax,n)}    & [\ref{ssub:method_pfct_path}] \\
  \bottomrule
  \end{tabular}
\egroup

%-------------------- details ---------------------------------------
\subsubsection{Constructor \tkzFct{pfct}{new(exprx,expry)}}
\label{ssub:method_pfct_new}

\texttt{Description: }Creates a parametric function object from two expressions
or callable Lua functions describing the components $x(t)$ and $y(t)$. When
strings are provided, they represent Lua expressions in the variable
\texttt{t}, evaluated using standard Lua rules (therefore requiring
\texttt{math.} prefixes).

\texttt{Arguments: }\par
\begin{itemize}
\item \code{exprx}: a string expression in \texttt{t}, or a callable Lua
      function representing $x(t)$;
\item \code{expry}: a string expression in \texttt{t}, or a callable Lua
      function representing $y(t)$.
\end{itemize}

\texttt{Returns: } a \tkzClass{pfct} object.

\subsubsection{Function \tkzFct{pfct}{compile(exprx,expry)}}
\label{ssub:method_pfct_compile}

\texttt{Description: }Compiles the two expressions defining $x(t)$ and $y(t)$
and returns the corresponding callable Lua functions (or wraps them into a
\tkzClass{pfct} object, depending on the implementation). This is a low-level
helper mainly intended for internal use.

\subsubsection{Method \tkzMeth{pfct}{x(t)}}
\label{ssub:method_pfct_x}

\texttt{Description: }Evaluates the $x$-component of the parametric curve at
parameter \code{t}.

\texttt{Returns: } a number.

\subsubsection{Method \tkzMeth{pfct}{y(t)}}
\label{ssub:method_pfct_y}

\texttt{Description: }Evaluates the $y$-component of the parametric curve at
parameter \code{t}.

\texttt{Returns: } a number.

\subsubsection{Method \tkzMeth{pfct}{point(t)}}
\label{ssub:method_pfct_point}

\texttt{Description: }Constructs the point $(x(t),y(t))$ associated with the
parametric curve. This method provides a direct bridge between numerical
evaluation and geometric construction.

\texttt{Returns: } a \tkzClass{point}.

\subsubsection{Method \tkzMeth{pfct}{path(tmin,tmax,n)}}
\label{ssub:method_pfct_path}

\texttt{Description: }Samples the parameter $t$ on the interval
$[tmin,tmax]$ using \code{n} subdivisions and returns the corresponding
\tkzClass{path}. Invalid values (NaN, infinities) may be skipped according to
the internal policy.

\texttt{Returns: } a \tkzClass{path}.

\texttt{Example: }
\begin{tkzexample}[latex=.4\textwidth]
\directlua{
 init_elements()
 PF.lis = pfct("sin(5*t)", "cos(3*t)")
 PA.curve = PF.lis:path(0, 2*math.pi, 400)
 }

\begin{center}
 \begin{tikzpicture}[scale=2]
  \tkzDrawCoordinates[smooth,cyan](PA.curve)
\end{tikzpicture}
 \end{center}
\end{tkzexample}


%--------------------------------------------------------------------
\subsection{Macros \tkzMacro{tkz-elements}{tkzDrawPointOnParamGraph} and
\tkzMacro{tkz-elements}{tkzDrawPointsOnParamGraph}}

These macros allow drawing points on a \emph{parametric curve} defined
in the module system.

The parametric function must already exist in the module table
\texttt{PF}.

%--------------------------------------------------------------------

\subsubsection{Macro \tkzMacro{tkz-elements}{tkzDrawPointOnParamGraph}}


\paragraph{Syntax}
\begin{verbatim}
\tkzDrawPointOnParamGraph[<TikZ options>]{<t>}{<name>}
\end{verbatim}

\paragraph{Description}
This macro draws a point belonging to a parametric curve stored
in the module \texttt{PF}.

The drawn point has coordinates
\[
\bigl(x(t),\,y(t)\bigr),
\]
where the parametric curve is defined by the object \texttt{PF.<name>}.

\paragraph{Arguments}
\begin{itemize}
  \item \texttt{<TikZ options>} (optional): graphical options applied to the point
  \item \texttt{<t>}: value of the parameter
  \item \texttt{<name>}: name of the parametric function stored in \texttt{PF}
\end{itemize}

\paragraph{Example}
\begin{verbatim}
\tkzDrawPointOnParamGraph[blue]{1.5}{lis}
\end{verbatim}

%--------------------------------------------------------------------
\subsubsection{Macro \tkzMacro{tkz-elements}{tkzDrawPointsOnParamGraph}}

\paragraph{Syntax}
\begin{verbatim}
\tkzDrawPointsOnParamGraph[<TikZ options>]{<t1,t2,...,tn>}{<name>}
\end{verbatim}

\paragraph{Description}
This macro draws several points belonging to a parametric curve stored
in the module \texttt{PF}.

For each value \(t_i\) in the list, a point of coordinates
\[
\bigl(x(t_i),\,y(t_i)\bigr)
\]
is computed and drawn.

\paragraph{Arguments}
\begin{itemize}
  \item \texttt{<TikZ options>} (optional): graphical options applied to all points
  \item \texttt{<t1,t2,...,tn>}: comma-separated list of parameter values
  \item \texttt{<name>}: name of the parametric function stored in \texttt{PF}
\end{itemize}

\paragraph{Example}
\begin{verbatim}
\tkzDrawPointsOnParamGraph[red]{0,0.5,1,1.5,2}{lis}
\end{verbatim}

\paragraph{Remarks}
\begin{itemize}
  \item These macros assume that the parametric function is already defined
        in the module \texttt{PF}.
  \item All evaluations are performed in Lua.
  \item The macros are compatible with parametric curves drawn using
        \texttt{\textbackslash tkzDrawCoordinates}.
\end{itemize}

\texttt{Example: }
\begin{tkzexample}[latex=.5\textwidth]
\directlua{
 init_elements()
 PF.lem = pfct("cos(t+pi/2)", "sin(2*t)")
 PA.courbe = PF.lem:path(0, 2*math.pi, 200)
 z.A = PF.lem:point(math.pi/2)
 }
\begin{center}
 \begin{tikzpicture}[scale=2]
\tkzGetNodes
  \tkzDrawCoordinates[smooth,cyan](PA.courbe)
  \tkzDrawPoint[blue](A)
  \tkzDrawPointsOnParamGraph[red]{0,2,3,5}{lem}
  \tkzDrawPointOnParamGraph[blue]{4}{lem}
\end{tikzpicture}
 \end{center}
\end{tkzexample}

\endinput
%--------------------------------------------------------------------
\input{TKZdoc-elements-metapost.tex}


\part{Mathematical and Computational Foundations}
 \section{Computational Model and Geometric Engine}
\subsection{Introduction}

The package \tkzNamePack{tkz-elements} is not merely a collection of geometric
construction commands. It implements a coherent computational geometry engine
based on a unified mathematical model.

This chapter presents the mathematical and computational foundations
underlying the package. While most users interact with high-level geometric
objects such as \tkzClass{point}, \tkzClass{line}, \tkzClass{circle},
\tkzClass{triangle}, or \tkzClass{conic}, all constructions ultimately rely
on a compact algebraic core and a carefully designed numerical strategy.

The purpose of this chapter is threefold:

\begin{itemize}
\item to describe the mathematical model used internally,
\item to explain the numerical robustness strategy,
\item to clarify the architectural principles governing the engine.
\end{itemize}

This chapter may be read independently as a technical overview of the
geometric engine.

\subsection{The Complex-Plane Model}

All planar points are represented internally as complex numbers:

\[
z = x + iy,
\]

where \(x\) and \(y\) are real coordinates.

This choice provides a compact and algebraically coherent framework for
geometric computations. Vector addition, subtraction, rotation,
and homothety become natural algebraic operations.

The class \tkzClass{point} therefore has a dual interpretation:
it behaves both as a geometric point in the Euclidean plane and as a complex
number in the algebraic model. This design avoids maintaining separate
vector and complex abstractions and provides a compact, coherent framework
for geometric computations.

The complex-plane representation allows:

\begin{itemize}
\item direct implementation of vector arithmetic,
\item concise expressions for rotations and symmetries,
\item natural encoding of orientation via complex multiplication,
\item simplified determinant and dot-product calculations.
\end{itemize}

The use of complex numbers ensures both elegance and computational efficiency.

%--------------------------------------------------------------------
\subsubsection{The Point Class as a Complex Number}
\label{sub:complex_numbers}

For example,
\begin{verbatim}
z.A = point(1, 2)
z.B = point(1, -1)
\end{verbatim}
define two point objects whose associated affixes are
\[
z_A = 1 + 2i,
\qquad
z_B = 1 - i.
\]
The notation \verb|z.A| refers to a Lua object stored in table \verb|z|,
whereas \(z_A\) denotes its associated complex number.

\medskip
If one prefers to work with standalone variables (without using reserved
tables), one may also write:
\begin{verbatim}
za = point(1, 2)
zb = point(1, -1)
\end{verbatim}
The only difference is organizational: \verb|z.A| is stored in the table
\verb|z|, while \verb|za| is a regular Lua variable.

\medskip
The algebraic structure is implemented through Lua metamethods and a small
set of class methods, allowing standard operators to perform geometric
operations in a natural way.

%--------------------------------------------------------------------
% Metamethods table
%--------------------------------------------------------------------
\begin{center}
  \bgroup
  \catcode`_=12
  \small
  \begin{minipage}{\textwidth}
  \captionof{table}{Point (complex) metamethods.}
  \begin{tabular}{lll}
    \toprule
    \textbf{Metamethod} & \textbf{Application} & \textbf{Result} \\
    \midrule
    \tkzMeta{point}{add(z1,z2)}    & \verb|z.a + z.b|        & affix \\
    \tkzMeta{point}{sub(z1,z2)}    & \verb|z.a - z.b|        & affix \\
    \tkzMeta{point}{unm(z)}        & \verb|-z.a|             & affix \\
    \tkzMeta{point}{mul(z1,z2)}    & \verb|z.a * z.b|        & affix \\
    \tkzMeta{point}{concat(z1,z2)} & \verb|z.a .. z.b|       & dot product (real number)\footnote{If \(O\) is the origin of the complex plane, then \(\,z_1 .. z_2\) corresponds to the dot product of the vectors \(\overrightarrow{Oz_1}\) and \(\overrightarrow{Oz_2}\).} \\
    \tkzMeta{point}{pow(z1,z2)}    & \verb|z.a ^ z.b|        & determinant (real number) \\
    \tkzMeta{point}{div(z1,z2)}    & \verb|z.a / z.b|        & affix \\
    \tkzMeta{point}{tostring(z)}   & \verb|tostring(z.a)|    & TeX-friendly display \\
    \tkzMeta{point}{tonumber(z)}   & \verb|tonumber(z.a)|    & affix or \texttt{nil} \\
    \tkzMeta{point}{eq(z1,z2)}     & \verb|z.a == z.b|       & boolean \\
    \bottomrule
  \end{tabular}
  \end{minipage}
  \egroup
\end{center}

%--------------------------------------------------------------------
% Methods table
%--------------------------------------------------------------------
\begin{center}
  \bgroup
  \catcode`_=12
  \small
  \begin{minipage}{\textwidth}
  \captionof{table}{Point (complex) class methods.}
  \begin{tabular}{lll}
    \toprule
    \textbf{Method} & \textbf{Application} & \textbf{Result} \\
    \midrule
    \tkzMeth{point}{conj()}  & \verb|z.a:conj()| & affix (conjugate) \\
    \tkzMeth{point}{mod()}   & \verb|z.a:mod()|  & real number (modulus) \\
    \tkzMeth{point}{abs()}   & \verb|z.a:abs()|  & real number (modulus) \\
    \tkzMeth{point}{norm()}  & \verb|z.a:norm()| & real number (squared modulus) \\
    \tkzMeth{point}{arg()}   & \verb|z.a:arg()|  & real number (argument in radians) \\
    \tkzMeth{point}{get()}   & \verb|z.a:get()|  & \verb|re| and \verb|im| (two real numbers) \\
    \tkzMeth{point}{sqrt()}  & \verb|z.a:sqrt()| & affix \\
    \bottomrule
  \end{tabular}
  \end{minipage}
  \egroup
\end{center}


\subsubsection{Example of complex use}

Let |za = math.cos(a) + i math.sin(a)| .
This is obtained from the library by writing

\begin{mybox}
   |za = point(math.cos(a),math.sin(a))|.
\end{mybox}

Then |z.B = z.A * za| describes a rotation of point A by an angle |a|.

\begin{minipage}{.45\textwidth}
\begin{verbatim}
\directlua{
 init_elements()
 z.O = point(0, 0)
 z.A = point(1, 2)
 a = math.pi / 6
 za = point(math.cos(a), math.sin(a))
 z.B = z.A * za}
\begin{tikzpicture}
 \tkzGetNodes
 \tkzDrawPoints(O,A,B)
 \tkzDrawArc[->,delta=0](O,A)(B)
 \tkzDrawSegments[dashed](O,A O,B)
 \tkzLabelAngle(A,O,B){$\pi/6$}
\end{tikzpicture}
  \end{verbatim}
\end{minipage}
\begin{minipage}{.55\textwidth}
\directlua{
 init_elements()
 z.O = point(0, 0)
 z.A = point(1, 2)
 a = math.pi / 6
 za = point(math.cos(a), math.sin(a))
 z.B = z.A * za}
\begin{center}
  \begin{tikzpicture}[scale=2]
  \tkzGetNodes
  \tkzDrawPoints(O,A,B)
  \tkzDrawArc[->,delta=0,thick](O,A)(B)
  \tkzDrawSegments[dashed](O,A O,B)
  \tkzLabelAngle(A,O,B){$\pi/6$}
  \end{tikzpicture}
\end{center}

\end{minipage}

\subsubsection{Point operations (complex)}

\begin{minipage}{.5\textwidth}
\begin{verbatim}
\directlua{
 init_elements()
 z.o = point(0, 0)
 z.a = point(1, -1)
 z.b = point(2, 1)
 z.bp = -z.b
 z.c = z.a + z.b
 z.d = z.a - z.b
 z.e = z.a * z.b
 z.f = z.a / z.b
 z.ap = point.conj(z.a)
 z.g = z.b * point(math.cos(math.pi / 2),
     math.sin(math.pi / 2))}

\begin{tikzpicture}
 \tkzGetNodes
 \tkzInit[xmin=-2,xmax=3,ymin=-2,ymax=3]
 \tkzGrid
 \tkzDrawSegments(o,a o,b o,c o,e o,b')
 \tkzDrawSegments(o,f o,g)
 \tkzDrawSegments[red](a,c b,c b',d a,d)
 \tkzDrawPoints(a,...,g,o,a',b')
 \tkzLabelPoints(o,a,b,c,d,e,f,g,a',b')
\end{tikzpicture}
\end{verbatim}
   \end{minipage}
\begin{minipage}{.5\textwidth}
\directlua{
 init_elements()
 z.o = point(0, 0)
 z.a = point(1, -1)
 z.b = point(2, 1)
 z.bp = -z.b
 z.c = z.a + z.b
 z.d = z.a - z.b
 z.e = z.a * z.b
 z.f = z.a / z.b
 z.ap = point.conj(z.a)
 z.g = z.b * point(math.cos(math.pi / 2), math.sin(math.pi / 2))}

\begin{center}
  \begin{tikzpicture}
           \tkzGetNodes
           \tkzInit[xmin=-2,xmax=3,ymin=-2,ymax=3]
            \tkzGrid
           \tkzDrawSegments(o,a o,b o,c o,e o,b' o,f o,g)
           \tkzDrawSegments[red](a,c b,c b',d a,d)
           \tkzDrawPoints(a,...,g,o,a',b')
           \tkzLabelPoints(o,a,b,c,d,e,f,g,a',b')
  \end{tikzpicture}
\end{center}

\end{minipage}

%--------------------------------------------------------------------
\subsubsection{Barycentric Combination}
\label{sub:barycentric_combination}

A fundamental operation in the computational engine is the
barycentric combination of points.

Given points \(z_1, \dots, z_n\) and associated real weights
\(w_1, \dots, w_n\), the barycenter is defined as

\[
G = \frac{\sum_{i=1}^{n} w_i z_i}{\sum_{i=1}^{n} w_i}.
\]

Because points are represented as complex numbers,
this operation reduces to a weighted linear combination.

\medskip
The internal function \code{barycenter\_} implements this formula
directly at the algebraic level. The public interface
\code{barycenter} provides a user-friendly wrapper.

\medskip
The barycentric construction is used extensively throughout
the package (in triangle centers, homotheties, affine constructions,
and other derived geometric algorithms).

\paragraph{Example.}


\begin{tkzexample}[latex=.35\textwidth]
\directlua{
 init_elements()
 z.A = point(1, 0)
 z.B = point(5, -1)
 z.C = point(2, 5)
 z.G = tkz.barycenter({z.A, 3}, {z.B, 1}, {z.C, 1})
}
\begin{center}
\begin{tikzpicture}
\tkzGetNodes
\tkzDrawPolygon(A,B,C)
\tkzDrawPoints(A,B,C,G)
\tkzLabelPoints(A,B,C,G)
\end{tikzpicture}
\end{center}
\end{tkzexample}


\subsection{Algebraic Primitives and Geometric Operators}

The computational engine relies on a minimal set of algebraic
operations defined on complex numbers.

These primitives are sufficient to implement all higher-level
geometric constructions such as projections, intersections,
membership tests, orientation checks, and affine combinations.

Because points are represented as complex numbers,
vector operations reduce to elementary algebra.

\subsubsection{Vector Arithmetic}

Given two points represented by complex numbers,
vector subtraction is defined naturally:

\[
\overrightarrow{AB} = z_B - z_A.
\]

Addition and scalar multiplication follow directly
from complex arithmetic.

These operations form the basis of translations,
midpoints, affine combinations, and homotheties.


\subsubsection{Dot Product}

Given two vectors represented by complex numbers
\(z_1 = a+ib\) and \(z_2 = c+id\),
their scalar product is defined as

\[
z_1 \verb|..| z_2 = ac + bd.
\]

In the implementation, this quantity is computed
using the operator \verb|..|:

\[
z_1 \verb|..| z_2 = ac + bd.
\]

Equivalently, the dot product corresponds to the real part of
the product of the first vector and the conjugate of the second:

\[
z_1 \verb|..| z_2
=
\operatorname{Re}\big(z_1 \,\overline{z_2}\big).
\]

\medskip
The dot product plays a central role in planar geometry.
It is used for:

\begin{itemize}
\item orthogonality tests,
\item projections onto a line,
\item distance computations,
\item angle measurements.
\end{itemize}

\paragraph{Distance computation.}

The squared distance between two points \(A\) and \(B\) is

\[
d(A,B)^2 =
(z_B - z_A) \verb|..|(z_B - z_A).
\]

\paragraph{Angle measurement.}

For three points \(A,B,C\), the angle \(\widehat{ABC}\) satisfies

\[
\cos(\theta)
=
\frac{
(z_A - z_B) \verb|..| (z_C - z_B)
}{
\lVert z_A - z_B \rVert
\,
\lVert z_C - z_B \rVert
}.
\]


\subsubsection{Determinant and Oriented Area}

Given two vectors represented by complex numbers
\(z_1 = a+ib\) and \(z_2 = c+id\),
their determinant is defined as

\[
\det(z_1, z_2) = ad - bc.
\]

This quantity represents the signed (or oriented) area
of the parallelogram generated by the two vectors.

In the implementation, the determinant is computed
using the operator \verb|^|:

\[
z_1 \verb|^| z_2 = ad - bc.
\]

\medskip
Given three points \(A,B,C\), the oriented area of triangle \(ABC\)
is obtained from

\[
\det(\overrightarrow{AB},\overrightarrow{AC})
=
(z_B - z_A) \verb|^| (z_C - z_A).
\]

The sign of this quantity determines the orientation:

\begin{itemize}
\item positive: direct (counterclockwise),
\item negative: indirect (clockwise),
\item zero (up to tolerance): aligned points.
\end{itemize}

This primitive is fundamental for:

\begin{itemize}
\item triangle orientation,
\item segment membership tests,
\item line intersection,
\item inside/outside classification.
\end{itemize}


\subsection{Numerical Robustness and Tolerance Control}


Geometric computations rely on floating-point arithmetic and are therefore
subject to numerical approximation errors.

To ensure robustness, all membership and position tests use a configurable
tolerance parameter:

\[
\texttt{tkz.epsilon}
\]

This tolerance is applied systematically in:

\begin{itemize}
\item alignment tests,
\item intersection detection,
\item tangency classification,
\item equality comparisons,
\item geometric membership tests.
\end{itemize}

Rather than relying on strict equality, the engine evaluates whether
a quantity is sufficiently close to zero.

This strategy guarantees numerical stability across complex constructions.

\subsection{Geometric Classification Model}


Geometric objects provide a unified classification interface based on
three possible states:

\begin{center}
\texttt{"ON"} \quad \texttt{"IN"} \quad \texttt{"OUT"}
\end{center}

This tri-state model applies consistently across:

\begin{itemize}
\item \tkzClass{line}
\item \tkzClass{circle}
\item \tkzClass{triangle}
\item \tkzClass{conic}
\end{itemize}

Membership and region tests are therefore harmonized throughout the system.

For backward compatibility, boolean interfaces are preserved where necessary,
but the internal model is uniformly tri-state.

\subsection{Degenerate Configurations}


Special care is taken to manage degenerate geometric cases:

\begin{itemize}
\item coincident points,
\item concentric circles,
\item aligned triangle vertices,
\item zero-radius constructions,
\item tangency limit cases.
\end{itemize}

Rather than allowing numerical instability, the engine classifies these
configurations explicitly.

This design choice prevents undefined behaviors and improves the reliability
of higher-level constructions.

\subsection{Architecture of the Computational Engine}

The engine is structured in two layers:

\subsection*{Primary Computational Functions}

Low-level functions perform algebraic computations.
They are internal and are not part of the public API.

\subsection*{Object-Oriented Layer}

Geometric classes provide user-facing methods.
These methods call primary functions while ensuring:

\begin{itemize}
\item consistent tolerance handling,
\item type safety,
\item backward compatibility.
\end{itemize}

This separation maintains clarity between computation and abstraction.

\subsubsection{Internal Data Tables}

In Lua, the main data structure is the \emph{table}. It functions both as an array and as a dictionary, allowing you to store sequences of values, associate keys with values, and even represent complex objects. Tables are the foundation for representing geometric structures such as points, lines, and triangles in \tkzNamePack{tkz-elements}.

\subsubsection{General Use of Tables}

Tables are the only data structure "container" integrated in Lua.
 They are associative arrays  which associates a key (reference or index) with a value in the form of a field (set) of key/value pairs. Moreover, tables have no fixed size and can grow based on our need dynamically.

Tables are created using curly braces:

\begin{mybox}
\begin{verbatim}
  T = {} % T is an empty table.
\end{verbatim}
\end{mybox}

In the next example, \code{coords} is a table containing four points written as coordinate strings. Lua indexes tables starting at 1.

\begin{tkzexample}[latex=0.35\textwidth]
\directlua{
 local coords = { "(0,0)", "(1,0)", "(1,1)", "(0,1)" }
 for i, pt in ipairs(coords) do
    tex.print(i .. ": " .. pt)
    tex.print([[\\]])
 end}
\end{tkzexample}

You can define a table with unordered indices as follows:

\begin{mybox}
\begin{verbatim}
coords = {[1] = "(0,0)", [3] = "(1,1)", [2] = "(1,0)"}
\end{verbatim}
\end{mybox}

Accessing an indexed element is straightforward:

\begin{verbatim}
   tex.print(coords[3]) --> (1,1)
\end{verbatim}

You can also append new elements:

\begin{verbatim}
coords[4] = "(0,1)"
\end{verbatim}

To iterate over a table with numeric indices in order, use \code{ipairs}:

\begin{mybox}
\begin{verbatim}
for i, v in ipairs(coords) do
print(i, v)
end
\end{verbatim}
\end{mybox}

This will print:

\begin{verbatim}
1 (0,0)
2 (1,0)
3 (1,1)
4 (0,1)
\end{verbatim}


\textbf{Key–Value Tables.}

Tables may also use string keys to associate names to values. This is especially useful in geometry for storing attributes like color or label:

\begin{mybox}
\begin{verbatim}
properties = {
A = "blue",
B = "green",
C = "red"
}
properties.D = "black"
\end{verbatim}
\end{mybox}

To iterate over all key–value pairs (order not guaranteed), use \code{pairs}:

\begin{mybox}
\begin{verbatim}
for k, v in pairs(properties) do
print(k, v)
end
\end{verbatim}
\end{mybox}

\textbf{Deleting an entry.}

To remove an entry, assign \code{nil} to the corresponding key:

\begin{verbatim}
properties.B = nil
\end{verbatim}


\subsubsection{Variable Number of Arguments}

Tables can also be used to store a variable number of function arguments:

\begin{mybox}
\begin{verbatim}
function ReturnTable(...)
return table.pack(...)
end
\end{verbatim}
\end{mybox}

\begin{mybox}
\begin{verbatim}
function ParamToTable(...)
local mytab = ReturnTable(...)
for i = 1, mytab.n do
print(mytab[i])
end
end

ParamToTable("A", "B", "C")
\end{verbatim}
\end{mybox}

This technique is useful for collecting points, coordinates, or any variable-length data.

\textbf{Accessing with Sugar Syntax.}

In tables with string keys, there are two common syntaxes:

\begin{itemize}
\item \code{properties["A"]} — always valid;
\item \code{properties.A} — shorter, but only works with string keys that are valid Lua identifiers (not numbers).
\end{itemize}

\subsubsection{Table \code{z}}

The most important table in \tkzNamePack{tkz-elements} is \code{z}, used to store geometric points. It is declared as:

\begin{verbatim}
z = {}
\end{verbatim}

Each point is then stored with a named key:

\begin{mybox}
\begin{verbatim}
z.M = point(3, 2)
\end{verbatim}
\end{mybox}

The object \code{z.M} holds real and imaginary parts, accessible as:

\begin{verbatim}
z.M.re --> 3
z.M.im --> 2
\end{verbatim}

If you print it:

\begin{verbatim}
tex.print(tostring(z.M))
\end{verbatim}

You get the complex representation \code{3+2i}.

Moreover, points in the \code{z} table are not just data—they are objects with methods. You can perform geometric operations directly:

\begin{verbatim}
z.N = z.M:rotation(math.pi / 2)
\end{verbatim}

This makes the \code{z} table central to object creation and manipulation in \tkzNamePack{tkz-elements}.



\subsection{Design Principles and Trade-offs}


The design of \tkzNamePack{tkz-elements} follows several guiding principles:

\begin{itemize}
\item Mathematical coherence,
\item Minimal algebraic core,
\item Numerical robustness,
\item Backward compatibility,
\item Performance efficiency.
\end{itemize}

The choice of complex representation, tri-state classification,
and tolerance-based comparisons results from balancing
mathematical rigor and practical usability.



\clearpage\newpage
\small\printindex

\end{document}